\documentclass[12pt,a4paper]{article}

\usepackage{fontspec}
\setmainfont{TeX Gyre Pagella}
\usepackage{microtype}
\usepackage{geometry}
\usepackage{setspace}
\usepackage{amsmath,amssymb,amsthm}
\usepackage{mathtools}
\usepackage[hidelinks]{hyperref}
\usepackage{xcolor}
\usepackage{booktabs}
\usepackage{array}
\usepackage{url}
\usepackage{xurl}

\geometry{
  a4paper,
  top=2.5cm,
  bottom=2.5cm,
  left=3cm,
  right=3cm
}

\onehalfspacing

\hypersetup{
  pdftitle={Constraints of Asymptotically Safe Quantum Gravity on
            Scalar-Tensor Theories and Cosmic Strings},
  pdfauthor={Flyxion}
}

\title{\textbf{Constraints of Asymptotically Safe Quantum Gravity\\
on Scalar-Tensor Theories and Cosmic Strings}}

\author{Flyxion}

\date{\today}

\begin{document}

\maketitle

\begin{abstract}
\noindent
This essay examines the constraints that the asymptotic-safety paradigm
imposes on two prominent classes of phenomenological models: shift-symmetric
Horndeski scalar-tensor theories proposed as explanations of dynamical dark
energy, and Abelian Higgs models capable of producing cosmic strings. Using
the Functional Renormalization Group (FRG) as the central analytical
instrument, we demonstrate that the requirement of ultraviolet (UV)
completeness through quantum scale symmetry severely restricts the viable
infrared (IR) parameter space of both model classes. In the Horndeski sector,
the cubic braiding coupling is shown to be symmetry-protected and vanishes
identically at the shifted Gaussian fixed point, rendering kinetic braiding
models incapable of generating the observed dark energy density. In the cosmic
string sector, the Abelian Higgs model cannot accommodate spontaneous symmetry
breaking at sub-Planckian scales without the introduction of at least
forty-two additional Dirac fermions. We place these results in the broader
context of the Hubble tension, the NANOGrav 15-year dataset, and the
conceptual architecture of the asymptotically safe landscape, arguing that
UV consistency acts as a structural filter that collapses the apparently large
space of beyond-the-Standard-Model candidates into a much smaller set of
genuinely viable theories.
\end{abstract}

\newpage

\tableofcontents

\newpage

%%==========================================================================
\section{Introduction: The Crisis in Cosmology and the Need for UV Completion}
%%==========================================================================

Contemporary theoretical cosmology finds itself at an unusual impasse. The
Standard Model of Cosmology, $\Lambda$-CDM, rests on two invisible pillars:
cold dark matter, which provides the gravitational scaffolding for large-scale
structure, and the cosmological constant $\Lambda$, a static energy density
intrinsic to empty space. For several decades this framework provided
excellent agreement with observation. As measurement precision has improved,
however, significant internal tensions have emerged that resist accommodation
within the $\Lambda$-CDM paradigm. The most comprehensive recent analysis of
cosmological parameters from the Planck satellite \cite{Planck2018} and the
independent distance-ladder measurements assembled by Riess and collaborators
\cite{Riess2019} are in tension at more than $4\sigma$, establishing what is
now widely called the Hubble tension as a genuine crisis rather than a
statistical fluctuation.

The most discussed manifestation of this crisis is the Hubble tension. Local
measurements of the expansion rate $H_0$ via Cepheid variable stars and Type Ia
supernovae consistently yield higher values than those inferred from the CMB
assuming $\Lambda$-CDM cosmology \cite{Riess2019,Planck2018}. A second tension
concerns the parameter $\sigma_8$, which characterizes the amplitude of matter
clustering on large scales. CMB-derived predictions of $\sigma_8$ exceed those
inferred from weak gravitational lensing surveys of the late-time universe,
suggesting that $\Lambda$-CDM may overpredict the rate of structure growth
\cite{Mukhanov2005,Dodelson2003}. Neither tension constitutes a definitive
falsification of $\Lambda$-CDM, but together they suggest the presence of
physics beyond the static cosmological constant.

This situation has motivated an enormous literature on dynamical dark energy
(DDE), much of it invoking scalar fields coupled to gravity through derivative
interactions \cite{Clifton2012}. The most mathematically systematic framework
for such theories is Horndeski gravity \cite{Horndeski1974}, the most general
scalar-tensor theory producing second-order equations of motion in four
spacetime dimensions. Within this class, kinetic braiding models are
particularly widely studied: the Galileon and its relatives \cite{Nicolis2009}
mix the kinetic term of the scalar field with spacetime curvature in a way
that can drive accelerated expansion without introducing a bare cosmological
constant, and they reproduce the correct speed of gravitational waves $c_t = 1$
when the non-minimal coupling functions are appropriately restricted.

What is typically absent from such constructions is an account of their
ultraviolet behavior. Phenomenological dark energy models are built as
effective field theories valid at cosmological scales of order $H_0^{-1}$,
with no guarantee of mathematical consistency at the Planck scale $M_{\rm Pl}
\approx 10^{19}\,\text{GeV}$. This is not merely an aesthetic deficiency. A
theory whose couplings diverge or become ill-defined at finite energy is not
a theory of nature; it is a placeholder for one \cite{Polchinski1984,
ZinnJustin2002}. The requirement that a viable model admit a well-defined,
interacting UV limit is therefore a genuine physical constraint, not a
matter of theoretical preference.

The asymptotic-safety program provides precisely this constraint. Proposed
originally by Weinberg \cite{Weinberg1979} and realized concretely through
the Functional Renormalization Group framework developed by Wetterich
\cite{Wetterich1993} and subsequently applied to gravity by Reuter
\cite{Reuter1998}, asymptotic safety posits that the renormalization group
flow of quantum gravity approaches a non-Gaussian fixed point (NGFP) in the
deep UV. At this fixed point, quantum fluctuations balance each other out in a
condition of quantum scale symmetry. The theory remains interacting but finite,
free of Landau poles and perturbative divergences. Crucially, the fixed-point
structure imposes rigid constraints on the matter sector that must coexist with
quantum gravity \cite{Eichhorn2019,EichhornSchiffer2022handbook}. Not every
low-energy effective theory can be connected by a smooth renormalization group
trajectory to an asymptotically safe UV fixed point. Those that cannot are, in
the terminology introduced by Vafa and elaborated by subsequent researchers,
swampland candidates---mathematically problematic extrapolations that cannot
represent fundamental physics \cite{EichhornHebeckerPawlowskiWalcher2024}.

The present essay synthesizes recent results establishing these constraints
for two model classes of direct phenomenological relevance. The first is the
family of shift-symmetric Horndeski theories with kinetic braiding, which
are among the most popular candidate explanations of dynamical dark energy
\cite{EichhornLinoWagner2023}. The second is the Abelian Higgs model, the
minimal field-theoretic mechanism for cosmic string production
\cite{EichhornLinoMiqueleto2023}. For each model class, we describe the FRG
methodology, state the main theoretical results, and connect them to
observational data, particularly the NANOGrav 15-year dataset
\cite{NANOGrav2023} and the Hubble tension measurements.

The conclusion that emerges is what we term a small-landscape picture: the
requirement of asymptotic safety collapses the apparently vast space of
beyond-Standard-Model phenomenological models into a very small set of
genuinely UV-consistent theories \cite{BuoninfanteEtal2025}. This is not a
limitation of the paradigm but its most powerful feature.

%%==========================================================================
\section{The Asymptotic-Safety Paradigm: Fixed Points, Scaling, and Predictivity}
%%==========================================================================

\subsection{The Ultraviolet Problem in Quantum Gravity}

General Relativity is perturbatively non-renormalizable. Power-counting
arguments show that Newton's coupling $G$ has mass dimension $-2$ in four
spacetime dimensions, so higher-loop diagrams introduce operators of
arbitrarily high dimension with independent coupling constants \cite{Peskin1995,
Weinberg1995}. To renormalize the theory order by order requires an infinite
number of counterterms, which means an infinite number of free parameters that
must be fixed by measurement. The theory loses predictive power at the Planck
scale and above. This is the ultraviolet problem of quantum gravity.

Two broad strategies have been developed to address it. The first invokes new
degrees of freedom at short distances---strings, loops, or higher-dimensional
structures---that alter the theory's short-distance behavior by changing its
fundamental constituents. The second, which is the approach underlying
asymptotic safety \cite{Weinberg1979}, remains within the framework of quantum
field theory and the metric, but seeks a non-trivial UV fixed point under the
renormalization group. This fixed point plays the role that asymptotic freedom
plays in non-Abelian gauge theories: it ensures that all physical couplings
remain finite as the energy scale $k$ is taken to infinity.

The key insight, due originally to Wilson and Kogut \cite{WilsonKogut1974,
Wilson1971} and formalized for gravity by Weinberg \cite{Weinberg1979}, is that
perturbative non-renormalizability does not preclude the existence of a
well-defined non-perturbative theory. A theory is asymptotically safe if the
renormalization group flow, when expressed in terms of dimensionless couplings,
has an interacting fixed point with a finite-dimensional UV-critical surface.
The dimensionality of this critical surface determines the number of
experimentally measurable parameters, and hence the predictive content of the
theory \cite{ReuterSaueressig2012,Percacci2017}.

\subsection{The Functional Renormalization Group and the Wetterich Equation}

The primary technical tool of the asymptotic-safety program is the Functional
Renormalization Group, specifically the exact flow equation for the
scale-dependent effective action $\Gamma_k$ derived by Wetterich
\cite{Wetterich1993} (see Appendix~\ref{app:wetterich} for a derivation):
\begin{equation}
k\partial_k\Gamma_k
= \frac{1}{2}\,\mathrm{Tr}
\!\left[\bigl(\Gamma^{(2)}_k + R_k\bigr)^{-1} k\partial_k R_k\right].
\label{eq:wetterich}
\end{equation}
Here $\Gamma^{(2)}_k$ denotes the second functional derivative of the
effective action with respect to the fields, evaluated at a background
configuration, and $R_k$ is an infrared regulator that suppresses the
contribution of momentum modes below the sliding scale $k$. As $k$ is lowered
from a high initial scale $\Lambda_{\rm UV}$ toward zero, progressively more
fluctuations are integrated out and $\Gamma_k$ flows from a bare UV action to
the full quantum effective action $\Gamma \equiv \Gamma_{k \to 0}$. Equation
(\ref{eq:wetterich}) has a one-loop structural form but is exact: no
approximations have been made in its derivation \cite{Berges2002,Dupuis2021}.
All approximations enter when one truncates the infinite-dimensional theory
space to a manageable finite-dimensional subspace---a physically motivated set
of field operators whose couplings are tracked. The quality of the truncation
determines the reliability of the resulting beta functions and fixed-point
analysis \cite{Pawlowski2007}.

A fixed point of the flow is a point $\{g^*_i\}$ in the space of
dimensionless couplings where all beta functions vanish:
\begin{equation}
\beta_i(g^*) \equiv k\partial_k g_i\big|_{g = g^*} = 0.
\end{equation}
The stability properties of the fixed point are encoded in the stability
matrix $B_{ij} = \partial\beta_i/\partial g_j|_{g^*}$, whose eigenvalues
$\theta_I$ are the critical exponents (see Appendix~\ref{app:stability}).
A coupling with $\text{Re}(\theta_I) > 0$ is \emph{relevant}: it can deviate
from its fixed-point value as the scale flows to the IR, and must be fixed by
experiment. A coupling with $\text{Re}(\theta_I) < 0$ is \emph{irrelevant}:
it is driven back to its fixed-point value and therefore constitutes a
prediction of the theory, not a free parameter. The number of relevant
directions equals the number of free parameters of the asymptotically safe
theory---a finite number, in sharp contrast to the infinite-parameter expansion
of perturbative quantum gravity \cite{Goldenfeld1992,Cardy1996}.

\subsection{The Non-Gaussian Fixed Point of Quantum Gravity}

Evidence for the existence of a non-Gaussian fixed point (NGFP) in the
Einstein-Hilbert truncation, where $\Gamma_k$ is restricted to the Ricci
scalar and cosmological constant, was established by Reuter \cite{Reuter1998}
and has since been confirmed in progressively larger truncations by Reuter and
Saueressig \cite{ReuterSaueressig2012} and many subsequent works reviewed by
Eichhorn \cite{Eichhorn2019}, Percacci \cite{Percacci2017}, and Eichhorn and
Schiffer \cite{EichhornSchiffer2022handbook}. The NGFP is characterized by
finite values of the dimensionless Newton coupling $g = Gk^2$ and the
dimensionless cosmological constant $\lambda = \Lambda/k^2$, together with
two relevant directions corresponding to deformations of these two parameters.
The fixed-point values and critical exponents are regulator-dependent but
qualitatively robust: there exists an interacting UV fixed point around which
the theory is well-defined and predictive.

The inclusion of matter fields modifies the fixed-point structure. Scalar
fields, fermions, and gauge bosons each contribute differently to the beta
functions for the gravitational couplings, and the gravitational fluctuations
in turn contribute to the beta functions for the matter couplings
\cite{Eichhorn2022status,EichhornSchiffer2022handbook}. It is this
gravitational backreaction on the matter sector that generates the constraints
on phenomenological models that are the primary subject of the present essay.

\subsection{The Weak-Gravity Bound}

A central organizing concept for matter-gravity systems in the
asymptotic-safety framework is the weak-gravity bound
\cite{EichhornSchiffer2022handbook}. When gravitational fluctuations are
included, interacting fixed points in the matter sector that exist in the
absence of gravity may be displaced or destabilized. As the dimensionless
Newton coupling $g$ increases from zero, interacting fixed points in the matter
sector move through theory space. Beyond a critical value $g_{\rm crit}$, pairs
of fixed points collide and annihilate into the complex plane, leaving no real
fixed point. The matter theory thereby loses its UV completion.

Quantitatively, within the Einstein-Hilbert truncation and for $\lambda = 0$,
this critical value is found to be approximately $g_{\rm crit} \approx 1.4$.
The precise number is truncation-dependent, but the qualitative mechanism is
robust: strong gravitational fluctuations destroy the fixed-point structure of
the matter sector \cite{EichhornSchiffer2022handbook,Eichhorn2019}. The
weak-gravity bound therefore defines a forbidden region in the gravitational
coupling plane and constrains both the values of the gravitational couplings at
the UV fixed point and the matter content of any UV-complete gravity-matter
system.

%%==========================================================================
\section{Theory Space and the Architecture of RG Flow}
%%==========================================================================

\subsection{Theory Space as a Dynamical System}

To formulate the asymptotic-safety condition with mathematical precision, it
is useful to introduce the notion of theory space. Theory space $\mathcal{T}$
is the space of all action functionals consistent with the field content and
symmetries under consideration \cite{WilsonKogut1974,Polchinski1984}. Each
point in $\mathcal{T}$ corresponds to a specific effective action
\begin{equation}
\Gamma = \sum_i g_i \mathcal{O}_i,
\end{equation}
where $\mathcal{O}_i$ denote operators constructed from the fields and their
derivatives, and $g_i$ are the associated coupling constants.

The renormalization group defines a flow on $\mathcal{T}$:
\begin{equation}
k\partial_k g_i = \beta_i(g_1,g_2,\ldots).
\end{equation}
The set of beta functions therefore defines a vector field on theory space,
and the RG flow corresponds to the integral curves of this vector field
\cite{Wilson1971,Cardy1996,Goldenfeld1992}. The RG flow can equally be
interpreted as a trajectory in the coupling space, moving from the UV regime
toward the IR as the resolution scale $k$ is lowered. This interpretation,
originating in the block-spin approach of Kadanoff \cite{Kadanoff1966} and
Wilson's subsequent formalization \cite{Wilson1971,WilsonKogut1974}, reveals
the renormalization group as a coarse-graining operation on the physical
degrees of freedom: each step in the flow integrates out a thin shell of
momentum modes and produces a new, equivalent description at a lower
resolution.

The UV-critical surface is the subset of $\mathcal{T}$ consisting of all
RG trajectories that reach the UV fixed point as $k \to \infty$. Its
codimension equals the number of irrelevant directions at the fixed point.
Physical theories are those whose RG trajectories lie on the UV-critical
surface, since only they possess a UV-complete definition. All other points
in $\mathcal{T}$ correspond to effective field theories that eventually lose
predictive power at the energy scale where the couplings diverge.

\subsection{Canonical Scaling and Dimensionless Couplings}

Renormalization group analysis is most naturally expressed in terms of
dimensionless couplings \cite{Peskin1995,Weinberg1995,ZinnJustin2002}.
If a coupling $g_i$ has canonical mass dimension $d_i$, we define the
dimensionless quantity
\begin{equation}
\tilde g_i = k^{-d_i} g_i.
\end{equation}
The RG equation for $\tilde g_i$ then takes the form
\begin{equation}
k\partial_k \tilde g_i = -d_i \tilde g_i + \beta_i^{\rm quantum}(\tilde g).
\label{eq:dimless_rg}
\end{equation}
The first term represents canonical scaling arising from dimensional analysis,
while the second term encodes quantum corrections generated by fluctuations
\cite{Berges2002,Dupuis2021}.

For the Horndeski truncation considered in this work, the cubic braiding
coupling $h$ has canonical dimension $-3$ and the quartic coupling $g$ has
canonical dimension $-4$ in four spacetime dimensions. Both interactions are
therefore power-counting irrelevant in perturbation theory, with canonical
scaling terms $+3\tilde h$ and $+4\tilde g$ in equation (\ref{eq:dimless_rg}).
Their fate in the ultraviolet is determined by the interplay between this
canonical suppression and quantum corrections from gravitational fluctuations
\cite{EichhornLinoWagner2023}. In the absence of gravity the canonical scaling
dominates and both couplings are irrelevant; the question is whether gravity
introduces sufficient antiscreening to promote either coupling to relevance.
The answer, as we shall see, is no.

\subsection{Information Geometry of Coupling Space}

The space of couplings $\{g_i\}$ itself carries a natural geometric structure
when equipped with a metric \cite{Fisher1922,AmariNagaoka2000,Amari2016}.
The Fisher information matrix
\begin{equation}
G_{ij}(\theta) = \mathbb{E}\!\left[\partial_i \log P\;\partial_j \log P\right]
\end{equation}
defines a Riemannian metric on the parameter space of a statistical model.
In the renormalization group context, this idea motivates an information-geometric
interpretation of RG trajectories as geodesics on the manifold of coupling
constants, where the metric measures the distinguishability between nearby
effective theories \cite{Ruppeiner1995}. Fixed points are then critical points
of an entropy-like functional on theory space, and the stability matrix
$B_{ij}$ encodes the local curvature of this landscape.

While the information-geometric perspective does not alter the technical content
of the FRG analysis, it provides a clarifying conceptual framework. The
asymptotically safe fixed point is a state of maximal self-consistency: the
effective theory at every scale is identical, up to a rescaling of units, to
the theory at every other scale. This is the field-theoretic realization of
scale symmetry, and the Fisher information geometry makes precise the sense in
which all UV-complete theories are parametrically close to one another in the
vicinity of the fixed point \cite{Amari2016}.

\subsection{Symmetry Constraints on Renormalization Group Flows}

Symmetries of the effective action impose powerful constraints on the structure
of beta functions, and this constraint is responsible for some of the most
robust predictions in the asymptotic-safety literature \cite{Eichhorn2019}.
Suppose the action is invariant under a discrete symmetry transformation
$\Phi \to \mathcal{S}\Phi$. Operators in the action may transform either
trivially or with a sign under this symmetry. If an operator $\mathcal{O}_i$
is odd under the symmetry while all other operators are even, then the
corresponding beta function must take the form
\begin{equation}
\beta_i = g_i\,f(g_j),
\end{equation}
for some function $f$ of the remaining couplings. The coupling $g_i$ therefore
factors out of its own beta function. Consequently, $g_i = 0$ is always a
fixed point of the RG flow. Such a coupling is \emph{symmetry-protected}:
quantum fluctuations cannot generate it radiatively if it vanishes in the
ultraviolet \cite{EichhornLinoWagner2023}.

This mechanism is far more general than the specific Horndeski application
considered here. It applies whenever a theory space contains operators of mixed
parity under a preserved symmetry, and its consequences are exact within any
truncation that respects the symmetry. The practical implication is that
symmetry analysis, rather than detailed numerical computation, can often
determine the fate of a coupling at the fixed point \cite{EichhornSchiffer2022handbook}.

%%==========================================================================
\section{Shift-Symmetric Horndeski Gravity and the Asymptotic-Safety Constraint}
%%==========================================================================

\subsection{Horndeski Theories and Kinetic Braiding}

Horndeski gravity \cite{Horndeski1974} identifies the most general scalar-tensor
theory in four dimensions whose equations of motion remain second-order in
derivatives. This property prevents the introduction of additional propagating
degrees of freedom with the wrong sign kinetic term (Ostrogradski ghosts),
making Horndeski theories the preferred framework for stable dynamical dark
energy models \cite{Clifton2012,Nicolis2009}. Within this landscape,
shift-symmetric theories occupy a special position. Shift symmetry is the
invariance of the action under the constant shift $\phi \to \phi + a$. This
symmetry restricts all interactions to involve only derivatives of $\phi$,
and the fundamental kinetic variable is
\begin{equation}
\chi = -\frac{1}{2}\nabla_\mu\phi\,\nabla^\mu\phi.
\end{equation}

The kinetic braiding sector studied here includes three operators beyond the
Einstein-Hilbert term: the standard kinetic term $\chi$, the cubic braiding
term $h\chi\Box\phi$, and a quartic interaction $g\chi^2$. The coupling $h$
parametrizes the braiding between the scalar field's kinetic term and the
metric, giving rise to the braiding parameter $\alpha_B$ which appears in
modified gravitational slip and deviations of the effective dark energy
equation of state from $-1$. Observational constraints from the simultaneous
detection of gravitational waves and their electromagnetic counterpart
GW170817 established that the speed of tensor modes must satisfy
$|c_t - 1| \lesssim 10^{-15}$ \cite{Abbott2016}, which forces $G_{4,X} = 0$
and $G_5 = 0$ within the Horndeski framework and leaves the cubic braiding
term as the primary candidate mechanism for kinetic dark energy.

\subsection{FRG Truncation for Horndeski Models}

To analyze the asymptotic-safety constraints on these models, one truncates
the effective action as
\begin{equation}
\Gamma_k = \int d^4x\,\sqrt{g}\left[-\frac{M_{\rm Pl}^2}{2}(R - 2\Lambda)
+ \chi + h\chi\Box\phi + g\chi^2 + \zeta\chi R\right].
\end{equation}
The gravitational sector is truncated to the Einstein-Hilbert form, with
dimensionless Newton coupling $G$ and dimensionless cosmological constant
$\Lambda$. The gauge-fixing is performed in the Landau-DeWitt gauge
$\alpha = 0$, $\beta = 0$, which simplifies the propagator structure and
enables analytic evaluation of momentum traces. The metric fluctuations are
decomposed into irreducible components using the York decomposition
(Appendix~\ref{app:york}), separating the spin-2 transverse-traceless modes,
vector modes, and scalar trace modes. The Litim regulator
\begin{equation}
R_k(p^2) = (k^2 - p^2)\,\Theta(k^2 - p^2)
\end{equation}
is chosen for its analytic tractability with the Heaviside step function
$\Theta$, which reduces momentum integrals in the trace to simple rational
functions of $k$ \cite{Berges2002,Dupuis2021}. The evaluation of functional
traces uses heat-kernel techniques summarized in Appendix~\ref{app:heatkernel}.

\subsection{The Shifted Gaussian Fixed Point and the Vanishing of Braiding}

The central result of the FRG analysis of shift-symmetric Horndeski models is
the identification of the shifted Gaussian fixed point (SGFP). This fixed point
inherits its name from the fact that, in the absence of gravity, the couplings
$h$ and $g$ both have a Gaussian fixed point at zero; the inclusion of
gravitational fluctuations shifts the fixed-point value of $g$ while leaving
that of $h$ at zero \cite{EichhornLinoWagner2023}.

The vanishing of $h$ at the fixed point is not accidental; it is a consequence
of the symmetry structure described in Section~3.4. The operator $\chi\Box\phi$
transforms with a definite sign under the discrete $\mathbb{Z}_2$ transformation
$\phi \to -\phi$ while $\chi$ is invariant: the cubic term breaks this kinetic
$\mathbb{Z}_2$ symmetry. Because the remaining operators in the truncation
preserve this symmetry, the beta function for $h$ must take the form
\begin{equation}
\beta_h = h\,f(G,\Lambda,g) + \mathcal{O}(h^3),
\end{equation}
for some function $f$ of the remaining couplings. This means that $h = 0$ is
always a solution of $\beta_h = 0$, independently of the values of the
gravitational couplings. Quantum gravity fluctuations cannot generate the
braiding coupling radiatively from zero. The SGFP therefore satisfies
$h_* = 0$ exactly \cite{EichhornLinoWagner2023}.

The critical exponent associated with $h$ at the SGFP is negative, confirming
that $h$ is an irrelevant coupling. Its infrared value is predicted by the RG
flow: it remains zero at all scales below the Planck scale, driven to its
fixed-point value by the RG evolution. The braiding parameter $\alpha_B$
vanishes identically. Even with the inclusion of the non-minimal coupling
$\zeta\chi R$, the sign of the critical exponent for $h$ does not change,
reinforcing the robustness of the conclusion \cite{EichhornLinoWagner2023,
EichhornSchiffer2022handbook}.

\subsection{Phenomenological Failure and the Magnitude Mismatch}

The prediction $h_* = 0$ creates a fundamental conflict with the
phenomenological requirements of dynamical dark energy. To replace the
cosmological constant with a dynamical scalar field whose energy density
matches the observed value $\rho_\Lambda \approx (2\times 10^{-3}\,\text{eV})^4$,
the coupling $g$ must take the value
\begin{equation}
g_{\rm required} \approx \frac{1}{H_0^2 M_{\rm Pl}^2}
\approx 9 \times 10^7\,\text{eV}^{-4}.
\end{equation}
The RG analysis within the EFT framework yields an upper bound on the
infrared value of $g$ of order
\begin{equation}
g_{\rm predicted} \lesssim \frac{H^2}{M_{\rm Pl}^2}
\approx 10^{-7}\,\text{eV}^{-4},
\end{equation}
a discrepancy of approximately fifty-seven orders of magnitude
\cite{EichhornLinoWagner2023}. When $h = 0$ is substituted into the expression
for the scalar field energy density in the kinetic braiding model, one obtains
\begin{equation}
\rho_\phi = -\frac{\dot\phi^2}{4} < 0,
\end{equation}
which is negative---contradicting the basic requirement of a positive total
energy density. The shift-symmetric kinetic braiding model, precisely in the
limit imposed by asymptotic safety, yields a scalar field that contributes
negatively to the energy budget of the universe.

\subsection{Effective Cosmological Dynamics}

To connect the RG analysis directly to cosmological observables, consider the
scalar-tensor system on a spatially flat Friedmann-Lema\^itre-Robertson-Walker
background:
\begin{equation}
ds^2 = -dt^2 + a(t)^2 d\vec{x}^2.
\end{equation}
The Friedmann equation takes the form
\begin{equation}
H^2 = \frac{1}{3M_{\rm Pl}^2}\left(\rho_m + \rho_\phi\right),
\end{equation}
where $H = \dot a/a$ is the Hubble parameter and $\rho_m$ is the matter
density \cite{Mukhanov2005,Weinberg2008}. For the shift-symmetric scalar sector,
the energy density is
\begin{equation}
\rho_\phi = \dot\phi^2 G_{2,\chi} - G_2 + 3H\dot\phi^3\,G_{3,\chi}.
\end{equation}
The term proportional to $G_{3,\chi}$ corresponds to kinetic braiding and
generates deviations from the cosmological constant equation of state. If the
RG flow enforces $h = 0$, this contribution vanishes identically, and the
scalar field reduces to a minimally coupled kinetic component with
$\rho_\phi = -\dot\phi^2/4$. The renormalization group therefore constrains
cosmological dynamics by fixing the coefficients entering the background
equations \cite{EichhornLinoWagner2023,Peebles1980}. The absence of the
braiding term implies that the scalar sector cannot produce the dynamical dark
energy behavior required to resolve the Hubble tension within the
shift-symmetric Horndeski framework.

\subsection{Implications for the Hubble Tension}

The Hubble tension has stimulated proposals for dynamical dark energy as a
possible resolution: if the equation of state of dark energy departed from
$w = -1$ in the recent past, the inferred value of $H_0$ from the CMB might
shift upward toward the locally measured value \cite{Planck2018,Riess2019}.
Shift-symmetric kinetic braiding models are natural candidates for such a
scenario because they can in principle generate $w \neq -1$ and a non-trivial
gravitational slip.

The asymptotic-safety constraint eliminates this possibility at the level of
fundamental theory. If the braiding coupling $h$ is zero at all energy scales,
the scalar field cannot produce the kinetic dark energy required to generate a
departure from $w = -1$ \cite{EichhornLinoWagner2023}. Any resolution of the
Hubble tension through shift-symmetric Horndeski dynamics must therefore invoke
either non-polynomial truncations of the scalar potential, extensions beyond
the shift-symmetric sector, or a different class of scalar-tensor theories
entirely.

%%==========================================================================
\section{Cosmic Strings in the Asymptotically Safe Landscape}
%%==========================================================================

\subsection{Cosmic Strings and the Abelian Higgs Model}

Cosmic strings are one-dimensional topological defects formed during phase
transitions in the early universe when a continuous symmetry is spontaneously
broken \cite{Peebles1980}. In the simplest realization they arise from the
Abelian Higgs model: a $U(1)$ gauge theory with a charged complex scalar
$\phi$ whose potential admits a circle of vacua at $|\phi| = v$. When the
universe cools through the phase transition temperature, the scalar field
acquires a non-zero vacuum expectation value (VEV), and topologically stable
string-like configurations are formed at the boundaries between regions that
select different phases. The tension of such strings is set by the
symmetry-breaking scale $v$, and observational constraints from pulsar timing
arrays bound $v \lesssim 4.6\times 10^{14}\,\text{GeV}$ \cite{NANOGrav2023}.
This is below the Planck scale by approximately four orders of magnitude, a
fact that will prove critical when the asymptotic-safety constraints are applied
\cite{EichhornLinoMiqueleto2023}.

\subsection{Obstacles to UV Completion of the Abelian Higgs Model}

The Abelian Higgs model in the presence of gravity presents two principal
obstacles to asymptotic-safety completion \cite{EichhornLinoMiqueleto2023}.
The first concerns the fixed-point structure of the coupled gauge-scalar system.
In the absence of gravity, the $U(1)$ gauge coupling has a Landau pole: it
grows without bound at high energies, indicating a breakdown of the theory in
the UV. Gravitational antiscreening can in principle counteract this tendency
and generate an interacting fixed point for the gauge coupling, but the
detailed analysis shows that the interplay between gravitational fluctuations
and the gauge-scalar sector does not generically produce a fixed point with a
perturbatively stable scalar potential.

The second obstacle is a scale mismatch. Even in the presence of a free
(Gaussian) fixed point for the gauge sector, the symmetry-breaking VEV is
driven by the RG flow to the Planck scale:
\begin{equation}
v \sim M_{\rm Pl} \approx 10^{19}\,\text{GeV}.
\end{equation}
This is in severe conflict with the observational bound of
$v \lesssim 4.6\times 10^{14}\,\text{GeV}$ from NANOGrav \cite{NANOGrav2023}.
The minimal Abelian Higgs model, when embedded in asymptotic safety, produces
cosmic strings whose tension is vastly too large to be consistent with PTA data.

\subsection{The Fermion Requirement: The Forty-Two Dirac Fermion Barrier}

To force the symmetry-breaking scale below the Planck mass, one may extend
the model by coupling $N_f$ Dirac fermions charged under the $U(1)$ gauge
symmetry. Fermions contribute to the running of the gauge coupling with a sign
opposite to that of gauge bosons: they screen the coupling and can reduce the
fixed-point value of the gauge coupling $g_*^2$ relative to the Yukawa
coupling $y_*^2$, thereby shifting the locus of spontaneous symmetry breaking
to lower energies \cite{EichhornLinoMiqueleto2023}.

The analysis of the stability conditions reveals two distinct thresholds.
For the scalar potential to admit any real solution for the mass parameter
$m^2_*$ below the Planck scale, one needs at least $N_f \geq 22$ Dirac
fermions. However, full stability across the cosmologically relevant range of
the dimensionless cosmological constant requires
\begin{equation}
N_f \geq 42.
\end{equation}
This result converts a qualitative incompatibility into a precise
particle-content requirement: the simplest field-theoretic model of cosmic
strings demands at least forty-two additional charged Dirac fermions beyond
the Standard Model merely to satisfy the UV consistency conditions imposed by
asymptotic safety. From a model-building perspective, this is an extremely
high price for a simple infrared phenomenon---what one might aptly call a
complexity tax \cite{EichhornLinoMiqueleto2023}.

\subsection{Alignment with NANOGrav}

The theoretical obstacles to accommodating cosmic strings in asymptotic safety
find striking observational support in the NANOGrav 15-year dataset
\cite{NANOGrav2023}. This dataset provides strong evidence for a stochastic
gravitational-wave background (SGWB) via the detection of quadrupolar
Hellings-Downs angular correlations between pulsars. The spectral properties
of the observed signal, however, are inconsistent with the characteristic
spectrum expected from a network of stable field-theoretic cosmic strings.
The data prefer instead a population of supermassive black-hole binary systems
in the final stages of inspiral.

This observational preference aligns with the theoretical expectation from
asymptotic safety \cite{EichhornLinoMiqueleto2023}: the minimal models that
would produce a detectable cosmic-string background are precisely those that
are theoretically inconsistent without a dramatic extension of the matter
sector. The convergence of top-down theoretical constraints and bottom-up
observational evidence strengthens the case that simple field-theoretic cosmic
strings are not a generic prediction of UV-complete theories, but a feature
of fine-tuned or theoretically problematic constructions.

%%==========================================================================
\section{The Small Landscape: UV Consistency as a Structural Filter}
%%==========================================================================

\subsection{Theory Space and the UV Critical Surface}

The results described in the preceding sections instantiate a general principle:
the requirement of asymptotic safety dramatically reduces the dimensionality
of the viable parameter space \cite{Weinberg1979,ReuterSaueressig2012}. The
full theory space $\mathcal{T}$ is infinite-dimensional. The asymptotically
safe fixed point defines a UV-critical surface: the set of all RG trajectories
that pass through or terminate at the fixed point as $k \to \infty$. Points
not on this surface cannot be connected to the fixed point and therefore do
not correspond to UV-complete theories.

The dimensionality of the UV-critical surface equals the number of relevant
directions at the fixed point \cite{Eichhorn2019,Percacci2017}. In the
Einstein-Hilbert truncation this is two, corresponding to the Newton coupling
and the cosmological constant. Including matter fields typically adds a small
number of additional relevant directions. The key point is that this
dimensionality is finite, while $\mathcal{T}$ is infinite-dimensional. The
UV-critical surface is therefore a measure-zero subset of theory space. The
failure of kinetic braiding and simple cosmic strings is not surprising; it is
the generic expectation, with the remarkable thing being if an arbitrarily
chosen low-energy model happened to lie on the UV-critical surface.

\subsection{Predictive Power and Irrelevant Couplings}

The mechanism through which asymptotic safety constrains phenomenology is the
irrelevance of couplings at the UV fixed point. An irrelevant coupling is not
physically unimportant: it is predicted. Its infrared value is determined by
the RG flow from the fixed point, not by a free parameter that must be measured
\cite{WilsonKogut1974,Berges2002}. When the predicted infrared value conflicts
with what is observationally required, the model is excluded.

This mechanism has been spectacularly successful in the gravitational sector.
The quartic Higgs self-coupling is irrelevant at the gravitational NGFP,
meaning that its value at the electroweak scale is predicted by the RG flow
from the Planck scale. This prediction yields $M_H \approx 126\,\text{GeV}$
\cite{ReuterSaueressig2012,deBritoEichhornPfeiffer2023a}, in remarkable
agreement with the measured value of $125.25\,\text{GeV}$. The braiding
coupling $h$ and the cosmologically required value of $g$ fail an analogous
test in the opposite direction: the predicted values are incompatible with
observation. The same framework has been applied to Standard Model couplings
beyond the Higgs sector, including the Yukawa coupling of the top quark
\cite{deBritoEichhornRay2023}, the gauge coupling of hypercharge
\cite{deBritoEichhornLino2022}, and the quark and lepton mixing angles
\cite{EichhornGyftopoulos2025}.

\subsection{A Unified Constraint Mechanism}

The two phenomenological examples examined in this essay---shift-symmetric
Horndeski dark energy and Abelian Higgs cosmic strings---illustrate a common
structural mechanism by which asymptotic safety constrains infrared physics.
In both cases the decisive ingredient is the classification of couplings as
relevant or irrelevant at the ultraviolet fixed point \cite{Eichhorn2019,
EichhornSchiffer2022handbook}.

For the Horndeski sector, the braiding coupling $h$ is symmetry-protected and
irrelevant, forcing its infrared value to vanish \cite{EichhornLinoWagner2023}.
For the Abelian Higgs model, the scalar potential parameters are driven toward
Planck-scale values unless additional fermionic matter is introduced to modify
the RG flow \cite{EichhornLinoMiqueleto2023}. These examples demonstrate that
the ultraviolet consistency requirement acts as a structural filter on theory
space \cite{BuoninfanteEtal2025}. Phenomenological models that appear viable
at low energies may fail when embedded in a UV-complete framework. The
asymptotic-safety paradigm therefore reduces the apparent freedom of model
building and narrows the landscape of consistent theories.

\subsection{Comparison with the Swampland Program}

The asymptotic-safety constraints described in this essay bear a structural
resemblance to the swampland conjectures of string theory \cite{EichhornHebeckerPawlowskiWalcher2024}.
Both programs impose top-down consistency conditions on effective field theories,
restricting the space of viable low-energy models. The swampland program
typically invokes the absence of global symmetries, the absence of stable
non-supersymmetric vacua, or the Trans-Planckian Censorship Conjecture to
argue that many apparently consistent EFTs cannot be embedded in a
UV-complete framework \cite{EichhornHebeckerPawlowskiWalcher2024}.

The asymptotic-safety filter operates differently but achieves a similar
result through a more concrete, computational mechanism. Rather than invoking
general consistency arguments about quantum gravity, it works through explicit
fixed-point analysis of the coupled gravity-matter FRG flow. Its predictions
are therefore more specific and, in principle, more falsifiable. The braiding
coupling is not merely argued to be small; it is computed to be exactly zero.
The fermion count required for cosmic string viability is not merely bounded
below; it is fixed to a specific integer threshold. This concreteness is one
of the most attractive features of the asymptotic-safety approach to the
swampland problem \cite{EichhornHebeckerPawlowskiWalcher2024,
EichhornSchiffer2022handbook}.

%%==========================================================================
\section{Methodology: The FRG Toolchain for Gravity-Matter Systems}
\label{sec:method}
%%==========================================================================

\subsection{Truncation Strategy}

The practical implementation of the FRG for gravity-matter systems begins with
the choice of truncation \cite{Berges2002,Pawlowski2007,Dupuis2021}. The
general principle is that operators with lower canonical dimension dominate
the flow near the Gaussian or near-Gaussian fixed point, so truncations
including all operators up to a given canonical dimension capture the most
important physics. In practice, the Einstein-Hilbert sector with Newton
coupling and cosmological constant is the universal starting point,
supplemented by matter kinetic terms and the specific interaction operators
relevant to the model under study.

For the Horndeski analysis, the truncation described in Section~4.2 includes
the full set of shift-symmetric operators up to the cubic level, with the
quartic sector included as a robustness check through the coupling $\zeta$.
The decomposition of metric fluctuations into transverse-traceless, trace, and
longitudinal modes uses the York decomposition (Appendix~\ref{app:york}),
which separates the spin-2 transverse-traceless modes carrying the primary
gravitational degrees of freedom from the spin-0 trace mode that couples to
the scalar field through the $\zeta\chi R$ operator.

\subsection{The $P^{-1}F$ Expansion}

The evaluation of the right-hand side of the Wetterich equation requires
computing the trace of a matrix involving the full inverse propagator
$\Gamma^{(2)}_k + R_k$. This is accomplished through the $P^{-1}F$ expansion
\cite{Pawlowski2007,Dupuis2021}:
\begin{equation}
(\Gamma^{(2)}_k + R_k)^{-1}
= (P + F)^{-1}
= P^{-1}\sum_{n=0}^\infty (-F P^{-1})^n,
\end{equation}
where $P = \Gamma^{(2)}_k + R_k|_{\Phi=0}$ is the propagator evaluated at
a background field configuration and $F = \Gamma^{(2)}_k - \Gamma^{(2)}_k|_{\Phi=0}$
is the fluctuation-field-dependent part. The trace of each term in this
expansion can be evaluated using heat-kernel methods detailed in
Appendix~\ref{app:heatkernel}, or by direct momentum integration with the
Litim regulator. The extraction of beta functions by projection onto specific
field monomials is described in Appendix~\ref{app:projection}. The
approximation of neglecting gravitational and scalar anomalous dimensions,
justified near the shifted Gaussian fixed point where canonical scaling
dominates, simplifies the computation while preserving the qualitative
structure of the results \cite{Pawlowski2007,Dupuis2021}.

\subsection{Truncation Robustness and Regulator Dependence}

All practical implementations of the FRG require truncation of the
infinite-dimensional theory space. The reliability of results must therefore
be assessed through robustness checks \cite{Berges2002,Dupuis2021}.

Regulator dependence can be examined by repeating the calculation with
different infrared cutoff functions $R_k$. Although numerical values of
fixed-point couplings may vary slightly, the existence of the fixed point and
the sign of the critical exponents should remain stable \cite{Berges2002}.
Gauge dependence is controlled through the background-field method and the
choice of gauge-fixing parameters. The Landau-DeWitt gauge used in this work
simplifies the propagator structure while preserving background diffeomorphism
invariance \cite{Reuter1998,Percacci2017}.

Truncation stability can be tested by enlarging the truncation to include
higher-order operators. Existing studies of asymptotically safe gravity-matter
systems indicate that the existence of the gravitational fixed point and the
irrelevance of higher-derivative scalar operators are stable features across
a wide class of truncations \cite{EichhornSchiffer2022handbook,Eichhorn2022status}.
The symmetry-protection argument for $h_* = 0$ is particularly robust: it
does not depend on the specific numerical values of coefficients in the
truncation, but only on the discrete $\mathbb{Z}_2$ symmetry that the cubic
operator breaks. Any truncation preserving this symmetry will find $h_* = 0$
\cite{EichhornLinoWagner2023}.

%%==========================================================================
\section{Observational Connections and the Narrowing Landscape}
%%==========================================================================

\subsection{NANOGrav and the Gravitational Wave Background}

The 15-year dataset of the NANOGrav pulsar timing array provides the most
significant recent observational input for the theoretical constraints described
in this essay \cite{NANOGrav2023}. The detection of a stochastic
gravitational-wave background via Hellings-Downs angular correlations confirms
that PTAs are now sensitive instruments for cosmological physics. The spectral
index of the observed signal is critical: a network of field-theoretic cosmic
strings in the scaling regime produces a background with a characteristic
spectral shape determined by the string tension and the loop distribution
function. The NANOGrav data prefer a steeper spectrum more consistent with a
population of supermassive black-hole binaries in the nanohertz frequency band.

This observational preference aligns with the theoretical expectation from
asymptotic safety \cite{EichhornLinoMiqueleto2023}. The minimal Abelian Higgs
model produces strings with tension at the Planck scale, far too large to be
consistent with PTA constraints. The convergence of these two independent lines
of evidence---one top-down from quantum gravity and one bottom-up from
observation---represents a powerful joint constraint on the cosmic string
hypothesis. It exemplifies the connection noted by Eichhorn, Lino dos Santos,
and Miqueleto \cite{EichhornLinoMiqueleto2023}: quantum gravity constraints
on the particle-physics models that can give rise to a stochastic GW background
become directly testable through PTA observations.

\subsection{Gravitational Wave Astronomy as a Probe of Quantum Gravity}

The connection between asymptotic safety and the gravitational wave background
represents a methodological advance of significant importance. It demonstrates
that constraints derived from UV consistency at the Planck scale can be
translated into predictions for astrophysical observables in the nanohertz
frequency band. The bridge is the RG flow itself: the UV fixed-point structure
constrains the symmetry-breaking scale of the Abelian Higgs model, which in
turn determines the string tension, which determines the amplitude and spectral
shape of the gravitational wave background \cite{EichhornLinoMiqueleto2023}.

This chain of reasoning suggests that gravitational wave observatories, from
PTAs to LIGO \cite{Abbott2016} to future space-based detectors, are not
merely instruments for studying astrophysical compact objects but probes of
the fundamental structure of quantum gravity. A precise measurement of the
stochastic background spectrum can, in principle, discriminate between
field-theoretic string models embedded in different UV completions and thereby
test the predictions of asymptotic safety at cosmological scales
\cite{BuoninfanteEtal2025,Maggiore2008}.

\subsection{Black Holes, Regular Geometries, and the Asymptotically Safe Landscape}

Beyond cosmic strings and dark energy, the asymptotic-safety program has made
recent progress on other observational frontiers. The possibility of regular
black holes---spacetimes without a central singularity---has been studied within
asymptotically safe gravity through RG-improved metrics \cite{EichhornFernandes2025,
CarballoRubioEtal2025}. Such regular geometries can in principle be distinguished
from classical black holes through deviations in the shadow morphology visible
to the Event Horizon Telescope \cite{EichhornGoldHeld2023}. The positivity
bounds derived from unitarity and causality have also been applied within
asymptotically safe gravity to constrain the parameter space of the effective
theory \cite{EichhornSchifferPedersen2025}, and constraints on ultralight dark
matter have been derived from the gravitational fixed-point structure
\cite{AssantEichhornKnorr2025}.

%%==========================================================================
\section{Beyond the Current Paradigm: Open Questions and Future Directions}
%%==========================================================================

\subsection{Non-Polynomial Dynamics and the Deep IR}

The most significant theoretical limitation of current FRG analyses is the
reliance on polynomial truncations for the scalar sector \cite{Berges2002,
Dupuis2021}. The full effective potential $V(\phi)$ is an infinite-dimensional
function, and truncating it to a polynomial in $\phi^2$ captures only the
local behavior near the expansion point. Non-polynomial fixed-point potentials,
of the type found in scalar field theories using the local potential
approximation of the FRG, can exhibit qualitatively different behavior,
including the possibility of spontaneous symmetry breaking at scales well below
the truncation scale \cite{Goldenfeld1992,Stanley1971}.

For the Horndeski dark energy problem, a non-polynomial treatment of the
kinetic sector would amount to considering the general function $G_2(\chi)$
in the Horndeski Lagrangian rather than its truncation to a quadratic polynomial.
Whether such a treatment can reconcile asymptotic safety with dynamical dark
energy is an open and important question. The technical challenge is substantial:
the flow equation for a non-polynomial function is a functional partial
differential equation in field space, and its analysis requires either
numerical methods or specific ans\"atze for the functional form of the
potential.

\subsection{Lorentzian Signature and Causal Structure}

All existing FRG analyses of asymptotically safe gravity are performed in
Euclidean signature, where the metric is positive-definite and the path
integral is well-defined as a statistical sum. The connection to physical
Lorentzian spacetimes requires an analytic continuation that is well-understood
for static backgrounds but becomes problematic for dynamical cosmological
spacetimes where the causal structure is essential. Recent work on spin-foam
models and tensor field theories has suggested possible approaches to Lorentzian
FRG formulations \cite{EichhornGurauGyftopoulos2025,CastroEichhornGurau2026},
but a complete implementation for the full gravity-matter system remains
outstanding. This is particularly relevant for cosmological applications, where
the predictions of the UV-fixed-point analysis must be connected to the
real-time evolution of the universe.

\subsection{Area-Metric Gravity and Extended Geometries}

A further direction that may extend the reach of asymptotic safety is the
consideration of gravitational theories based on more general geometric
structures than the metric. Area-metric gravity, in which the fundamental
variable is an area metric $G_{\mu\nu\rho\sigma}$ rather than the standard
metric $g_{\mu\nu}$, provides more degrees of freedom and may admit richer
fixed-point structures \cite{BorissovaDittrichEichhornSchiffer2025}.
Scalar-Gauss-Bonnet gravity with higher-order couplings \cite{EichhornFernandes2026}
extends the Horndeski sector in directions not covered by the shift-symmetric
truncation analyzed here. These extensions may provide the additional structure
needed to reconcile asymptotic safety with specific phenomenological requirements,
or they may impose even more stringent constraints, advancing the understanding
of the asymptotically safe landscape in either case.

\subsection{From Proton Decay to Neutrino Masses: Matter-Gravity Unification}

The scope of asymptotic safety predictions extends well beyond the models
discussed in detail here. Recent work has established that the framework makes
predictions for proton stability \cite{EichhornRay2024}, neutrino mass
generation \cite{deBritoEichhornPereiraYamada2025}, and the possible existence
of axion-like particles \cite{deBritoEichhornLino2022}. The question of CPT
symmetry restoration has been addressed within the framework
\cite{EichhornSchifferCPT2025}, and higher-curvature operators in causal set
quantum gravity have been studied as a complementary UV-completion strategy
\cite{deBritoEichhornPfeiffer2023b}. The initial distribution of Standard Model
couplings at the Planck scale has been constrained by the structure of
asymptotically safe trajectories \cite{EichhornHeld2025}.

The breadth of these applications suggests that asymptotic safety is evolving
from a specific hypothesis about the UV behavior of quantum gravity into a
comprehensive framework for connecting Planck-scale physics to phenomenology
across all sectors of the Standard Model and its extensions
\cite{BuoninfanteEtal2025,Eichhorn2022status}. Whether this program will
ultimately yield a unique theory or a family of consistent completions depends
on questions about the global structure of theory space---the existence and
classification of universality classes of gravity-matter systems
\cite{CastroEichhornGurau2026}---that are not yet settled.

%%==========================================================================
\section{Diagrammatic Structure of Graviton-Matter Loop Contributions}
%%==========================================================================

\subsection{One-Loop Contributions to Matter Beta Functions}

The beta functions for matter couplings in an asymptotically safe gravity-matter
system receive contributions from two distinct sources: self-interactions of the
matter fields, which reproduce the flat-spacetime running, and gravitational
fluctuations mediated by virtual graviton exchange. Understanding the
diagrammatic structure of these contributions clarifies why gravity can promote
irrelevant couplings toward relevance and why the symmetry-protection argument
for the braiding coupling is particularly robust \cite{Reuter1998,Berges2002,
Dupuis2021}.

Within the one-loop nonperturbative structure of the Wetterich equation, the
trace on the right-hand side of equation~(\ref{eq:wetterich}) generates, after
expanding in the $P^{-1}F$ framework, a set of contributions that can be
organized diagrammatically. Consider a generic matter coupling $g_i$ appearing
as a vertex in the action. The gravitational contribution to $\beta_{g_i}$ arises
from diagrams in which two matter field lines are connected by a graviton
propagator. Schematically, one encounters two classes of diagram.

The first class involves a single graviton loop closing between two matter
vertices. This contribution is proportional to $g_i \cdot G_N / (k^2 + m_{\rm grav}^2)$,
where $G_N$ is the Newton coupling and $m_{\rm grav}^2 = 2\Lambda k^2$ is the
effective graviton mass generated by the cosmological constant \cite{Percacci2017,
ReuterSaueressig2012}. This term modifies the anomalous dimension of the matter
field and shifts the canonical scaling of $g_i$. For couplings with negative
canonical dimension $d_i < 0$, gravity can in principle compensate the
power-counting suppression and make the coupling marginal or relevant.

The second class involves the seagull diagrams generated by the minimal coupling
of matter to gravity through the covariant derivative and the metric determinant
$\sqrt{g}$. These diagrams produce contact interactions between the matter
sector and the graviton propagator without introducing additional matter vertices.
They contribute to the gravitational shift of fixed-point values but do not
alter the symmetry-based selection rules on which operators can be generated
\cite{EichhornSchiffer2022handbook}.

\subsection{The Gravitational Shift of Scalar Couplings}

For the scalar sector of the Horndeski truncation, the gravitational contributions
to the beta functions of $h$ and $g$ can be isolated by writing
\begin{align}
\beta_h &= \beta_h^{\rm matter} + \beta_h^{\rm grav}, \\
\beta_g &= \beta_g^{\rm matter} + \beta_g^{\rm grav}.
\end{align}
The matter contributions $\beta_h^{\rm matter}$ and $\beta_g^{\rm matter}$
reproduce the flat-spacetime running of these couplings in the absence of
gravity. In the shift-symmetric sector, $\beta_h^{\rm matter} \propto h$
by the symmetry argument of Section~3.4, and $\beta_g^{\rm matter}$ receives
contributions from the $\chi^2$ self-interaction and from the mixing of $h$
and $g$ through intermediate field configurations.

The gravitational contribution $\beta_h^{\rm grav}$ also inherits the
$\mathbb{Z}_2$ symmetry constraint: because gravity couples to the scalar
field only through $\chi = -\tfrac{1}{2}\nabla_\mu\phi\nabla^\mu\phi$, which
is $\mathbb{Z}_2$-even, any graviton-mediated loop diagram connecting to the
$h$-vertex must produce a factor proportional to $h$. The gravitational sector
therefore cannot break the $\mathbb{Z}_2$ symmetry protection of the braiding
coupling \cite{EichhornLinoWagner2023,EichhornSchiffer2022handbook}. This is
the diagrammatic expression of the symmetry argument: gravity mediates
interactions, but it cannot generate new quantum numbers that the matter sector
does not already carry.

\subsection{Antiscreening from Graviton Exchange}

The gravitational contribution to $\beta_g$ is qualitatively different: because
$g$ is $\mathbb{Z}_2$-even, the graviton loop can generate a non-zero
contribution to $\beta_g$ that does not vanish at $g = 0$. This contribution
takes the form
\begin{equation}
\beta_g^{\rm grav} \supset \frac{c_g\,G}{(1 - 2\Lambda)^p},
\end{equation}
where $c_g$ is a positive numerical coefficient determined by the graviton spin
content and $p$ depends on the regulator scheme \cite{Percacci2017,
EichhornLinoWagner2023}. This antiscreening contribution shifts the fixed-point
value of $g$ away from zero---the Gaussian fixed point---to a non-zero value
$g_* \neq 0$. The result is the shifted Gaussian fixed point: a fixed point in
which $h_* = 0$ is maintained by symmetry while $g_*$ is shifted to a positive
value by gravitational antiscreening.

The magnitude of this shift is controlled by the gravitational coupling strength
$G$ at the fixed point. Near the weak-gravity bound, where $G \to G_{\rm crit}$,
the shift can become large enough to change the sign of the critical exponent
$\theta_g$, potentially promoting $g$ from an irrelevant to a marginal or relevant
coupling. However, as shown by the detailed analysis of \cite{EichhornLinoWagner2023},
the resulting IR value of $g$ remains many orders of magnitude below the
phenomenologically required value even in this near-critical regime.

%%==========================================================================
\section{Operator Dimension Counting and the Truncation Landscape}
%%==========================================================================

\subsection{Canonical Dimensions and Power Counting}

A systematic understanding of which operators can appear at a given order in the
FRG truncation requires careful power counting \cite{WilsonKogut1974,Polchinski1984,
Peskin1995}. In four spacetime dimensions, an operator $\mathcal{O}$ with
canonical mass dimension $[\mathcal{O}] = d$ has a coupling constant of canonical
dimension $4 - d$ (from the requirement that the action be dimensionless). Operators
with $d > 4$ are power-counting irrelevant and their couplings have negative mass
dimension; operators with $d < 4$ are relevant; operators with $d = 4$ are marginal.

Table~\ref{tab:operators} lists the canonical dimensions and classical scaling
properties of the operators appearing in the Horndeski truncation studied in
this work, together with an indication of whether asymptotic safety predicts each
coupling to be relevant or irrelevant at the SGFP.

\begin{table}[h]
\centering
\renewcommand{\arraystretch}{1.4}
\begin{tabular}{llccc}
\toprule
Operator & Symbol & Canonical dim.\ of coupling & Classical scaling & AS verdict \\
\midrule
Cosmological constant & $\Lambda$ & $+2$ & Relevant & Free param. \\
Einstein-Hilbert & $R$ & $0$ & Marginal & Free param. \\
Scalar kinetic & $\chi$ & $0$ & Marginal & Free param. \\
Non-minimal coupling & $\zeta\chi R$ & $0$ & Marginal & Shifted, irrel. \\
Cubic braiding & $h\chi\Box\phi$ & $-3$ & Irrelevant & $h_* = 0$ (sym.) \\
Quartic interaction & $g\chi^2$ & $-4$ & Irrelevant & $g_* \neq 0$ (shifted) \\
Higher kinetic & $\chi^3$ & $-8$ & Strongly irrel. & $\approx 0$ \\
\bottomrule
\end{tabular}
\caption{Canonical dimensions and asymptotic-safety predictions for operators
in the shift-symmetric Horndeski truncation. The column ``AS verdict'' indicates
whether the coupling is a free parameter (relevant), symmetry-protected to zero,
or shifted to a small non-zero value by gravitational antiscreening.
Strongly irrelevant operators are suppressed by additional powers of
$(k/M_{\rm Pl})$ and are negligible in the near-Planck regime.}
\label{tab:operators}
\end{table}

\subsection{Beyond the Truncation: Convergence Arguments}

The operators listed in Table~\ref{tab:operators} constitute the minimal
truncation required to capture the phenomenology of kinetic braiding dark energy.
The question of whether extending this truncation to include higher-order operators
would qualitatively alter the conclusions can be addressed through several
arguments \cite{Berges2002,Dupuis2021,Pawlowski2007}.

First, the symmetry-protection argument for $h_* = 0$ is exact within any
truncation that preserves the $\mathbb{Z}_2$ symmetry. Adding higher-order
$\mathbb{Z}_2$-odd operators to the truncation would produce additional
symmetry-protected couplings, all vanishing at the fixed point, but would not
change the fate of $h$ itself. Second, the magnitude mismatch for $g$---fifty-seven
orders of magnitude between the AS prediction and the cosmologically required
value---is so large that it is not plausibly an artifact of the truncation.
Even if the non-polynomial treatment of $G_2(\chi)$ shifted $g_*$ by several
orders of magnitude, it would not bridge a gap of $10^{57}$ \cite{EichhornLinoWagner2023}.

Third, the convergence of FRG truncations in the scalar sector has been studied
extensively in the context of $O(N)$ models and related theories \cite{Berges2002,
Stanley1971,Goldenfeld1992}. In these cases the local potential approximation,
which retains the full non-polynomial potential $V(\phi)$ while truncating the
derivative sector, provides excellent quantitative agreement with exact results.
Applying this lesson to the Horndeski system suggests that extending the kinetic
truncation to the full function $G_2(\chi)$ is the most important next step,
while the qualitative predictions of the polynomial analysis remain reliable
at the level of order-of-magnitude estimates.

%%==========================================================================
\section{The RG Phase Diagram: From the Planck Scale to the Cosmic Horizon}
%%==========================================================================

\subsection{Structure of the Gravitational Phase Diagram}

The renormalization group flow of asymptotically safe gravity-matter systems
can be visualized as a phase diagram in the space of dimensionless couplings
$(G, \Lambda)$, supplemented by the matter coupling directions
\cite{ReuterSaueressig2012,Percacci2017,Eichhorn2019}. The phase diagram
organizes all RG trajectories and makes transparent which initial conditions
in the UV lead to viable low-energy physics.

In the Einstein-Hilbert sector alone, the phase diagram in the $(G, \Lambda)$
plane exhibits several qualitatively distinct regions. The non-Gaussian fixed
point (NGFP) sits at a specific location $(G_*, \Lambda_*)$ with $G_* > 0$ and
$\Lambda_* > 0$, determined by the condition $\beta_G = \beta_\Lambda = 0$. The
UV-critical surface is a two-dimensional manifold passing through this fixed
point, consisting of all RG trajectories that approach it as $k \to \infty$.
Trajectories that do not lie on this surface diverge from the NGFP and either
reach a singularity at finite scale or flow to a Gaussian fixed point at
$(G, \Lambda) = (0, 0)$.

Among the UV-critical trajectories, one must further identify those that
produce physically acceptable IR behavior: a positive effective Newton constant,
a small positive cosmological constant consistent with observation, and
matter sector couplings that match the known parameters of the Standard Model.
This additional requirement selects a subset of the UV-critical surface,
reducing the viable parameter space to what has been called the ``trajectory
of the universe'' in the asymptotic-safety literature \cite{ReuterSaueressig2012,
Bonanno2002}.

\subsection{Matter Coupling Trajectories}

When matter fields are included, the phase diagram extends into a higher-dimensional
space. The gravitational coupling directions remain $(G, \Lambda)$, while the
matter sector introduces additional dimensions for each independent coupling
\cite{Eichhorn2019,EichhornSchiffer2022handbook}. The structure of the RG flow
in this extended space determines the fate of each matter coupling.

For the Horndeski system, the relevant part of the phase diagram can be
described in terms of the three-dimensional space $(G, h, g)$ at fixed
$\Lambda = 0$ for simplicity. The SGFP is located at $(G_*, 0, g_*)$. The
UV-critical surface passes through this point with a two-dimensional tangent
plane spanned by the directions corresponding to $G$ and the linear combination
of $h$ and $g$ that is relevant at the fixed point. Since both $h$ and $g$ are
irrelevant at the SGFP, the UV-critical surface is actually two-dimensional (the
$G$ and $\Lambda$ directions), and trajectories that start with $h \neq 0$ in the
UV are driven toward $h = 0$ as $k \to \infty$ \cite{EichhornLinoWagner2023}.

Flowing from the UV fixed point toward the IR, trajectories initially move along
the two-dimensional UV-critical surface in the $(G, \Lambda)$ plane. As the scale
$k$ decreases through the Planck scale $M_{\rm Pl}$, the dimensionless couplings
begin to deviate from their fixed-point values. The Newton coupling $G = G_N k^2$
decreases as $k^2$, reproducing Newton's law at macroscopic scales. The
cosmological constant $\Lambda = \bar\Lambda/k^2$ grows at low scales, which in
the dimensionful theory corresponds to a small but non-zero value of $\bar\Lambda$
consistent with observation \cite{Bonanno2002,Planck2018}.

\subsection{The Separation of Scales and the Dark Energy Problem}

The most striking feature of the RG phase diagram for gravity-matter systems
is the enormous separation of scales between the UV fixed point at $k \sim M_{\rm Pl}$
and the infrared regime relevant for dark energy at $k \sim H_0$. This separation
spans approximately sixty orders of magnitude in energy:
\begin{equation}
\frac{M_{\rm Pl}}{H_0} \approx \frac{10^{19}\,\text{GeV}}{10^{-42}\,\text{GeV}}
= 10^{61}.
\end{equation}
During this enormous RG evolution, irrelevant couplings are suppressed by
powers of $(k/M_{\rm Pl})^{|\theta_I|}$. For a coupling with critical exponent
$|\theta_I| \sim 1$, the suppression factor at the scale $H_0$ is
$(H_0/M_{\rm Pl}) \approx 10^{-61}$. This is the origin of the fifty-seven-order-of-magnitude
mismatch for $g$: the canonical dimension of $g$ is $-4$, so the suppression
at the cosmological scale is $(H_0/M_{\rm Pl})^4 \sim 10^{-244}$ from canonical
scaling alone, partially offset by the gravitational shift of the fixed-point
value \cite{EichhornLinoWagner2023}.

This separation of scales is not a problem to be solved but a structural feature
of the theory that explains why the cosmological constant problem is so severe
from the EFT perspective: any mass scale in the theory is naturally of order
$M_{\rm Pl}$, and bringing it down to $H_0$ requires either enormous fine-tuning
or a symmetry mechanism. The asymptotic-safety prediction is that no such
mechanism exists within the shift-symmetric Horndeski framework. This is not
a failure of the paradigm; it is the paradigm successfully ruling out a class
of theories that could not have worked \cite{EichhornHebeckerPawlowskiWalcher2024,
BuoninfanteEtal2025}.

\subsection{Entropy, Information Loss, and Geometric Flows}

The structure of the RG phase diagram carries a deeper geometric interpretation
that connects the asymptotic-safety program to foundational work in geometric
analysis and nonequilibrium thermodynamics \cite{Fisher1922,AmariNagaoka2000,
Ruppeiner1995}. The renormalization group flow is, in a precise sense, an
information-loss process: as the scale $k$ is lowered, short-distance fluctuations
are integrated out and the effective description of the system becomes coarser.
The information content of the UV theory is progressively compressed into a
smaller number of relevant parameters \cite{Amari2016,Goldenfeld1992}.

This observation connects the FRG program to the mathematical theory of
geometric flows initiated by Hamilton's work on Ricci flow \cite{Hamilton1982}
and subsequently developed by Perelman's proof of the geometrization conjecture
\cite{Perelman2002}. In Ricci flow the metric of a Riemannian manifold evolves
according to
\begin{equation}
\partial_t g_{\mu\nu} = -2R_{\mu\nu},
\end{equation}
driving the geometry toward spaces of constant curvature---the fixed points of
the flow. The analogy with the RG is precise: just as Ricci flow eliminates
geometric irregularities to produce canonical geometries, the RG flow eliminates
short-scale fluctuations to produce scale-invariant effective theories. The
fixed points of Ricci flow (Einstein manifolds) correspond to fixed points of
the RG (scale-invariant quantum field theories) \cite{Gromov1981,Burago2001}.

Perelman's key insight was to introduce the $\mathcal{F}$-functional
\begin{equation}
\mathcal{F}(g,f) = \int_M (R + |\nabla f|^2)\,e^{-f}\,dV,
\end{equation}
which is monotonically non-decreasing along the Ricci flow and serves as a
Lyapunov functional ensuring convergence \cite{Perelman2002}. The analogous
object in the RG context is Zamolodchikov's $c$-function in two dimensions,
or more generally the proposed $a$-function in four dimensions, which decreases
monotonically along RG flows \cite{Cardy1996,ZinnJustin2002}. Whether a
monotone Lyapunov functional exists for the full gravity-matter RG flow of
asymptotic safety remains an open question, but the geometric analogy suggests
that such a functional, if it exists, would take the form of an entropy-like
quantity defined on theory space \cite{Jaynes1957,Onsager1931}.

The nonequilibrium thermodynamic perspective, pioneered by Onsager
\cite{Onsager1931} and later systematized by Prigogine \cite{Prigogine1967}
and developed into an information-theoretic framework by Jaynes \cite{Jaynes1957},
interprets irreversible processes as entropy-generating flows toward states of
maximum entropy. The RG flow, viewed as an irreversible coarse-graining
process, fits naturally into this framework: information is lost monotonically
as the resolution scale decreases, and fixed points are states of maximum
entropy compatible with the symmetry constraints of the theory
\cite{Amari2016,AmariNagaoka2000}. This interpretation does not alter the
technical content of the FRG analysis but suggests a broader conceptual
architecture in which asymptotic safety, geometric flow theory, and
nonequilibrium thermodynamics are unified aspects of a single information-geometric
framework \cite{Ruppeiner1995,Fisher1922}.

%%==========================================================================
\section{Consistency Checks and Comparison with Other Quantum Gravity Constraints}
%%==========================================================================

The asymptotic-safety constraints discussed in this work are part of a broader
effort to identify general principles that restrict the space of low-energy
effective field theories compatible with quantum gravity. Several independent
approaches have been developed in recent years, including the string-theoretic
swampland program \cite{EichhornHebeckerPawlowskiWalcher2024}, effective field
theory positivity bounds \cite{EichhornSchifferPedersen2025}, and semiclassical
consistency conditions. Although these frameworks arise from different
theoretical starting points, they share a common goal: to distinguish effective
theories that can arise from a consistent ultraviolet completion from those that
cannot. In this section we briefly compare the constraints derived from
asymptotic safety with those obtained from these alternative approaches.

\subsection{Relation to the Swampland Program}

The swampland program \cite{EichhornHebeckerPawlowskiWalcher2024}, originally
proposed in the context of string theory, seeks to characterize the set of
effective field theories that cannot be embedded into a consistent theory of
quantum gravity. Theories that fail these criteria are said to lie in the
swampland rather than in the landscape of viable quantum gravity vacua. Many
swampland conjectures place constraints on scalar-field potentials, derivative
interactions, and the range of validity of effective field theories. For
example, the de Sitter conjecture suggests that scalar potentials compatible
with quantum gravity must satisfy bounds of the form $|\nabla V| \geq c\,V$,
where $c$ is an $\mathcal{O}(1)$ constant in Planck units.

Although the asymptotic-safety framework does not assume string theory or extra
dimensions, it produces conceptually similar restrictions through an entirely
different mechanism \cite{Eichhorn2019,BuoninfanteEtal2025}. Couplings that are
irrelevant at the ultraviolet fixed point are not free parameters but predictions
of the theory. If the predicted infrared value of such a coupling conflicts with
the requirements of a phenomenological model, the model is excluded. In this
sense, asymptotic safety provides an explicit dynamical mechanism by which
certain low-energy theories are removed from the space of UV-complete models,
complementing the more axiomatic constraints of the swampland program
\cite{EichhornHebeckerPawlowskiWalcher2024}.

\subsection{Positivity Bounds from Effective Field Theory}

Another independent set of constraints arises from analyticity and unitarity
of scattering amplitudes \cite{EichhornSchifferPedersen2025}. In local,
Lorentz-invariant quantum field theories, dispersion relations imply positivity
conditions on the coefficients of higher-derivative operators in the low-energy
effective action. For scalar field theories the coefficient of the operator
$(\partial_\mu\phi\,\partial^\mu\phi)^2$ must be non-negative to ensure that
forward scattering amplitudes obey the optical theorem. Recent work has extended
these bounds to gravity-matter systems within asymptotically safe gravity
\cite{EichhornSchifferPedersen2025}. These analyses show that positivity bounds
can be compatible with the existence of a non-Gaussian fixed point but may
restrict the allowed range of higher-derivative couplings along RG trajectories
emerging from the fixed point. In this sense, positivity bounds provide
complementary infrared consistency checks on the RG flows predicted by
asymptotic safety.

\subsection{Semiclassical Gravity Constraints}

A third class of consistency conditions arises from semiclassical gravity.
In semiclassical treatments, quantum matter fields propagate on a classical
gravitational background, and their stress-energy tensor acts as a source in
the Einstein equations. Semiclassical consistency requires that expectation
values of the stress-energy tensor satisfy certain energy conditions and
stability criteria. The asymptotic-safety framework automatically incorporates
many of these constraints because the RG flow controls the sign and magnitude
of kinetic and interaction terms in the effective action \cite{Percacci2017,
ReuterSaueressig2012}. Operators that would lead to pathological behavior---such
as wrong-sign kinetic terms or unstable potentials---typically correspond to
directions in theory space that do not lie on the UV-critical surface. The
recent work on regular black holes and gravastars \cite{EichhornFernandes2025,
CarballoRubioEtal2025} represents a concrete application of this principle:
the RG-improved geometry removes the central singularity precisely because the
running Newton coupling suppresses gravitational strength at short distances,
maintaining semiclassical consistency across the interior.

\subsection{Convergence of Independent Constraints}

Taken together, these different approaches suggest a coherent picture.
Independent theoretical principles---renormalization group fixed points,
string-theoretic consistency conditions, analyticity of scattering amplitudes,
and semiclassical stability---tend to restrict effective field theories in
qualitatively similar ways \cite{BuoninfanteEtal2025,EichhornSchiffer2022handbook}.
The asymptotic-safety paradigm provides a concrete realization of this idea by
explicitly constructing renormalization group flows that connect ultraviolet
fixed points to infrared physics. The examples studied in this essay illustrate
how this connection translates Planck-scale consistency conditions into
observable constraints on cosmological models, and the convergence of these
independent lines of argument strengthens the case that the asymptotically safe
landscape is genuinely small.

%%==========================================================================
\section{Renormalization Group Flow as Coarse-Graining of Physics}
%%==========================================================================

The renormalization group can be understood intuitively as a systematic
procedure for coarse-graining physical systems \cite{WilsonKogut1974,Kadanoff1966,
Goldenfeld1992}. At any given scale $k$, the effective action $\Gamma_k$
describes the dynamics of fluctuations with momenta smaller than $k$.
Fluctuations with higher momentum have already been integrated out. Lowering
the scale $k$ corresponds to gradually averaging over shorter-distance physics.
In this sense, the renormalization group provides a continuous interpolation
between microscopic and macroscopic descriptions of nature.

This coarse-graining picture was originally developed in the context of
statistical physics, where it explains the universality of critical phenomena
\cite{Wilson1971,Kadanoff1966,Stanley1971,Cardy1996}. Systems with very
different microscopic structures can exhibit identical large-scale behavior
because the renormalization group flow drives their couplings toward the same
fixed point. The universality classes of critical phenomena---characterized by
the fixed-point structure and critical exponents---are among the most striking
examples of emergent simplicity from underlying complexity.

The asymptotic-safety scenario applies this logic to quantum gravity
\cite{Weinberg1979,Reuter1998}. Instead of describing the infrared behavior of
condensed matter systems, the renormalization group flow is followed toward
extremely high energies. If the flow approaches a fixed point as
$k \to \infty$, the theory remains well-defined at arbitrarily short distances.
Ultraviolet completion is not achieved by introducing new degrees of freedom
but by the self-organizing structure of the renormalization group flow itself,
which is precisely the appeal of the program from a minimalist perspective
\cite{Eichhorn2019,Percacci2017}.

%%==========================================================================
\section{Why Fixed Points Generate Predictivity}
%%==========================================================================

The predictive power of an asymptotically safe theory arises directly from the
structure of the ultraviolet fixed point. Near a fixed point, many couplings
are attracted toward specific values as the energy scale increases
\cite{WilsonKogut1974,Berges2002}. These couplings are called irrelevant because
their ultraviolet behavior is not sensitive to their initial values: any nearby
trajectory is drawn toward the fixed-point value. From the perspective of
low-energy physics, this means that the infrared values of irrelevant couplings
are determined by the RG trajectory emerging from the fixed point rather than
by independent free parameters.

A useful analogy comes from dynamical systems theory. A stable fixed point in
a phase portrait attracts trajectories from a neighborhood of initial conditions.
Even if the initial values of the couplings differ, they are drawn toward the
same fixed-point value as the energy scale increases. In asymptotically safe
quantum gravity, the UV fixed point plays precisely this role: only a small
number of directions in theory space allow trajectories to leave the fixed point
toward lower energies, corresponding to the relevant couplings that must be
measured experimentally. All other couplings are predicted \cite{Eichhorn2019,
ReuterSaueressig2012}. The irrelevance of the Horndeski braiding coupling $h$
is the specific realization of this mechanism in the dark energy sector
\cite{EichhornLinoWagner2023}.

%%==========================================================================
\section{Quantum Gravity as a Structural Constraint on Cosmological Models}
%%==========================================================================

The examples discussed in this essay illustrate a broader principle: cosmological
model building cannot be separated from the ultraviolet structure of fundamental
physics. In traditional effective field theory approaches, cosmological models
are constructed by writing down the most general Lagrangian consistent with
symmetries at the relevant energy scale. Parameters are adjusted to reproduce
observational data. This strategy has proven extremely successful, but it
implicitly assumes that the effective theory can be extended to arbitrarily high
energies without encountering inconsistencies \cite{Polchinski1984,ZinnJustin2002}.

Quantum gravity challenges this assumption. The renormalization group structure
of asymptotically safe gravity implies that not every low-energy effective theory
can be connected to a consistent ultraviolet completion
\cite{Weinberg1979,Reuter1998,Eichhorn2019}. Only those theories lying on RG
trajectories that originate from the ultraviolet fixed point are admissible.
From this perspective, ultraviolet consistency acts as a structural constraint
on cosmological model building: instead of beginning with the most general
Lagrangian allowed by symmetry, one must restrict attention to theories whose
parameters can arise from renormalization group flows compatible with the
fixed-point structure of quantum gravity.

The two case studies analyzed in this work illustrate how such constraints
operate in practice. In the shift-symmetric Horndeski sector, the renormalization
group flow forces the braiding coupling to vanish because it is symmetry
protected at the ultraviolet fixed point \cite{EichhornLinoWagner2023}. In the
cosmic string scenario, the Abelian Higgs model admits spontaneous symmetry
breaking only at the Planck scale when embedded in the asymptotically safe
framework, and lower symmetry-breaking scales require the introduction of a
large number of additional fermionic degrees of freedom
\cite{EichhornLinoMiqueleto2023}. Both examples reveal the same underlying
mechanism: the ultraviolet fixed point determines which directions in theory
space are physically accessible, and couplings that lie outside this set cannot
be realized by any RG trajectory connected to the fixed point.

%%==========================================================================
\section{Implications for the Strategy of Cosmological Model Building}
%%==========================================================================

The analysis presented in this work highlights a shift in the strategy of
theoretical cosmology. For many years, the dominant approach to model building
has been bottom-up: one begins with observational anomalies or unexplained
phenomena, proposes new fields or interactions capable of addressing them, and
investigates whether the resulting model is compatible with existing data
\cite{Clifton2012,Dodelson2003,Weinberg2008}. While this methodology has
produced a rich landscape of candidate theories, it often postpones questions
about ultraviolet consistency.

The asymptotic-safety framework reverses this order of reasoning
\cite{Weinberg1979,Eichhorn2019,EichhornSchiffer2022handbook}. Instead of
beginning with phenomenology and asking how it might be embedded into a
fundamental theory, one begins with the requirement that the theory remain
well-defined at arbitrarily high energies. This requirement restricts the
renormalization group trajectories that are physically admissible. Cosmological
models must then be constructed within the subset of effective theories
compatible with those trajectories.

From this perspective, ultraviolet completion becomes a primary criterion for
theoretical viability rather than a secondary consideration. The case studies
explored in this essay illustrate this shift concretely. The shift-symmetric
Horndeski models that have been widely studied as candidates for dynamical dark
energy rely on a coupling that is driven to zero by the ultraviolet fixed point
\cite{EichhornLinoWagner2023}. The minimal Abelian Higgs model predicts cosmic
strings with tensions far above the observational bounds once it is embedded
into the asymptotically safe framework \cite{EichhornLinoMiqueleto2023}. These
results do not imply that dynamical dark energy or cosmic strings are impossible;
they indicate that their realization within a UV-complete theory requires
additional structure absent from the simplest models. The renormalization group
therefore acts as a guide for model building by indicating which directions in
theory space are compatible with a fundamental completion, and by doing so it
connects the abstract requirement of quantum consistency to the concrete
observational program of precision cosmology \cite{Planck2018,NANOGrav2023,
BuoninfanteEtal2025}.

%%==========================================================================
\section{Beta Functions of the Abelian Higgs Model under Gravity}
\label{sec:abelian_beta}
%%==========================================================================

\subsection{Setup and Field Content}

The Abelian Higgs model coupled to quantum gravity provides the simplest
arena in which the competition between gravitational antiscreening and
matter-sector dynamics can be studied with full quantitative control. The
minimal field content consists of a $U(1)$ gauge field $A_\mu$, a complex
charged scalar $\phi$, and the metric $g_{\mu\nu}$ in the Einstein-Hilbert
truncation. The FRG effective action is \cite{EichhornLinoMiqueleto2023}
\begin{equation}
\Gamma_k = \int d^4x\sqrt{g}\left[
-\frac{M_{\rm Pl}^2}{2}(R - 2\Lambda)
+ \frac{1}{4}F_{\mu\nu}F^{\mu\nu}
+ |D_\mu\phi|^2 + m^2|\phi|^2 + \lambda_4|\phi|^4
\right],
\end{equation}
where $D_\mu = \partial_\mu - ig_{\rm g}A_\mu$ is the gauge-covariant
derivative. In terms of dimensionless couplings $g = G k^2$, $\lambda = \Lambda/k^2$,
and rescaled matter couplings $\tilde m^2 = m^2/k^2$, $\tilde\lambda_4 = \lambda_4$,
$\tilde g_{\rm g}^2 = g_{\rm g}^2/k^2$, the Wetterich equation generates beta
functions that govern the RG flow of all five parameters.

\subsection{Gauge Coupling Beta Function}

The running of the gauge coupling receives contributions from the gauge boson
and scalar loops (matter sector) and from the graviton loop (gravitational
sector). After projecting the flow equation onto the $F_{\mu\nu}F^{\mu\nu}$
operator and evaluating the functional traces with the Litim regulator, the
beta function for the dimensionless gauge coupling squared $\tilde g_{\rm g}^2$
takes the form \cite{EichhornLinoMiqueleto2023,deBritoEichhornLino2022}
\begin{equation}
\beta_{\tilde g_{\rm g}^2}
= \tilde g_{\rm g}^2\left[
-2
+ \frac{5 G}{18\pi(1-2\Lambda)^2}
- \frac{5G}{9\pi(1-2\Lambda)}
+ \frac{1}{48\pi^2}\,\tilde g_{\rm g}^2
\right].
\label{eq:beta_gauge}
\end{equation}
The term $-2$ arises from canonical scaling ($[\tilde g_{\rm g}^2] = -2$
in four dimensions). The terms proportional to $G$ represent gravitational
antiscreening: they are positive and grow as $\Lambda \to 1/2$, providing
the mechanism that can save the theory from the Landau pole. The final term
$\propto \tilde g_{\rm g}^4$ is the standard matter-sector screening contribution
from the charged scalar loop.

\subsection{Fixed-Point Topology and the Landau-Pole Resolution}

Setting $\beta_{\tilde g_{\rm g}^2} = 0$ and ignoring the matter sector term
for simplicity, the non-trivial fixed point is located at
\begin{equation}
\tilde g_{\rm g,*}^2
= \frac{(1+\tilde m^2)^2}{\sqrt{3}}\sqrt{40\pi G},
\end{equation}
where the scalar mass contribution enters through the propagator structure of
the heat-kernel trace. This fixed point exists for any $G > 0$, confirming
that gravitational antiscreening resolves the Abelian triviality problem: the
theory no longer hits a Landau pole but instead approaches a finite interacting
fixed point \cite{deBritoEichhornLino2022,EichhornLinoMiqueleto2023}.

The topology of the fixed-point landscape in the $(G, \tilde g_{\rm g}^2)$ plane
is therefore qualitatively different from the flat-spacetime case. In the
absence of gravity, the only fixed point of the gauge coupling is the Gaussian
one at $\tilde g_{\rm g,*}^2 = 0$, which is UV-attractive but corresponds to a
trivial, non-interacting theory. With gravity, an interacting fixed point
emerges at finite $\tilde g_{\rm g,*}^2 \propto \sqrt{G}$, providing a UV
completion for the interacting Abelian gauge theory.

\subsection{Scalar Sector Beta Functions}

The beta functions for the scalar mass and quartic coupling are
\begin{align}
\beta_{\tilde m^2}
&= \left(-2 + \eta_\phi\right)\tilde m^2
+ \frac{3\tilde\lambda_4}{16\pi^2(1+\tilde m^2)^2}
- \frac{3\tilde g_{\rm g}^2}{16\pi^2(1+\tilde m^2)}
+ f_m(G,\Lambda),
\label{eq:beta_mass}
\\
\beta_{\tilde\lambda_4}
&= \left(\eta_\phi\right)\tilde\lambda_4
- \frac{9\tilde\lambda_4^2}{16\pi^2(1+\tilde m^2)^3}
+ \frac{\tilde g_{\rm g}^4}{4\pi^2(1+\tilde m^2)^2}
+ f_\lambda(G,\Lambda),
\label{eq:beta_quartic}
\end{align}
where $\eta_\phi$ is the scalar anomalous dimension and $f_m$, $f_\lambda$
encode the gravitational contributions from the transverse-traceless and trace
scalar graviton modes. The gauge coupling appears in equation (\ref{eq:beta_mass})
with a negative sign, reflecting the antiscreening of the mass by gauge boson
fluctuations, and in equation (\ref{eq:beta_quartic}) with a positive sign,
reflecting the quartic coupling generated by box diagrams at one loop.

The central difficulty for cosmic string production is that the gravitational
contribution $f_\lambda(G,\Lambda)$ tends to drive $\tilde\lambda_{4,*}$
negative, destabilizing the scalar potential. Spontaneous symmetry breaking
requires a negative mass squared $\tilde m^2_* < 0$ at the fixed point
simultaneously with a positive quartic coupling $\tilde\lambda_{4,*} > 0$
to ensure boundedness from below. The minimal Abelian Higgs model coupled to
gravity generically fails to satisfy both conditions simultaneously
\cite{EichhornLinoMiqueleto2023}.

%%==========================================================================
\section{The Yukawa-Abelian-Higgs System and Fermionic Screening}
\label{sec:yukawa}
%%==========================================================================

\subsection{Extension to the Yukawa Sector}

The obstruction identified in Section~\ref{sec:abelian_beta} motivates the
inclusion of charged Dirac fermions $\psi_i$ ($i = 1,\ldots,N_f$) coupled to
the scalar through Yukawa interactions. The interaction Lagrangian takes the
form \cite{EichhornLinoMiqueleto2023}
\begin{equation}
\mathcal{L}_{\rm Yuk}
= i\bar y_k\left(\phi^*\bar\psi^c\psi + \phi\bar\psi\psi^c\right),
\end{equation}
where $\psi^c = C\bar\psi^T$ denotes the charge-conjugate spinor and $\bar y_k$
is the dimensionful Yukawa coupling. In terms of the dimensionless Yukawa
coupling $\tilde y^2 = \bar y^2/k^2$, the extended system has beta functions
for all couplings modified by fermionic loop contributions.

\subsection{Modified Gauge Coupling Beta Function}

With $N_f$ Dirac fermions, the gauge coupling beta function becomes
\begin{equation}
\beta_{\tilde g_{\rm g}^2}
= \tilde g_{\rm g}^2\left[
-2
+ \frac{5G}{18\pi(1-2\Lambda)^2}
- \frac{5G}{9\pi(1-2\Lambda)}
+ \frac{1+N_f}{48\pi^2}\,\tilde g_{\rm g}^2
\right].
\label{eq:beta_gauge_Nf}
\end{equation}
The fermion contribution enters the matter-sector term with coefficient
$(1+N_f)$: the scalar contributes $1$ and each Dirac fermion contributes
$1$ to the screening of the gauge coupling. The fixed-point value of the
gauge coupling therefore decreases as $N_f$ increases:
\begin{equation}
\tilde g_{\rm g,*}^2 \propto \frac{1}{1+N_f}.
\label{eq:gauge_fp_Nf}
\end{equation}
This screening is the mechanism by which fermions shift the locus of
spontaneous symmetry breaking to lower energies. A smaller $\tilde g_{\rm g,*}^2$
means that the gauge boson contribution to the scalar mass beta function
(\ref{eq:beta_mass}) is reduced, allowing $\tilde m^2_* < 0$ to be achieved
at a sub-Planckian scale.

\subsection{Yukawa Coupling Beta Function}

The Yukawa coupling itself evolves under the RG according to
\begin{equation}
\beta_{\tilde y^2}
= -2\tilde y^2
+ \frac{(4+N_f)\tilde y^4}{8\pi^2}
- \frac{3\tilde g_{\rm g}^2\,\tilde y^2}{8\pi^2}
+ f_y(G,\Lambda),
\end{equation}
where the canonical scaling term $-2\tilde y^2$ reflects the mass dimension $-1$
of the Yukawa coupling in four dimensions, and $f_y(G,\Lambda)$ contains
gravitational contributions. A non-trivial fixed point $\tilde y_*^2 > 0$
exists when the balance between the positive $(4+N_f)\tilde y^4$ term and the
gauge-coupling screening determines the fixed-point value.

The ratio $\tilde g_{\rm g,*}^2 / \tilde y_*^2$ plays the critical role in
determining whether the scalar potential admits spontaneous symmetry breaking.
Specifically, the fixed-point condition $\tilde m^2_* < 0$ is satisfied only
when $\tilde g_{\rm g,*}^2 / \tilde y_*^2$ is sufficiently small, which requires
$N_f$ large enough to screen the gauge coupling below the Yukawa coupling
\cite{EichhornLinoMiqueleto2023}.

\subsection{Quantitative Stability Thresholds}

The detailed stability analysis of the coupled fixed-point equations reveals
two distinct thresholds:

For $N_f < 22$, no parameter region in the $(G,\Lambda)$ plane admits
simultaneous solutions $\tilde m^2_* < 0$ and $\tilde\lambda_{4,*} > 0$.
The scalar potential is universally unstable: gauge boson fluctuations drive
the quartic coupling negative before the mass can become negative, producing
a potential unbounded from below.

For $22 \leq N_f < 42$, positive solutions for $\tilde\lambda_{4,*}$ exist but
only within a restricted ``strip'' in the $(\Lambda, G)$ plane at $\Lambda < 0$.
The allowed region is narrow and disappears as $\Lambda \to 0$, indicating that
cosmologically viable symmetry breaking requires a large negative cosmological
constant at the Planck scale.

For $N_f \geq 42$, the stability conditions are satisfied across a broad range
of the cosmological constant including a neighborhood of $\Lambda = 0$. The
cosmological constant sensitivity becomes $|\Lambda_{\rm threshold}| \lesssim 7.2$
for $N_f = 42$, improving as $N_f$ increases \cite{EichhornLinoMiqueleto2023}.
This threshold at $N_f = 42$ constitutes the complexity tax quantified in Section~5.3.

%%==========================================================================
\section{Explicit Beta Functions for Shift-Symmetric Horndeski Couplings}
\label{sec:horndeski_beta}
%%==========================================================================

\subsection{The Complete Beta Functions at Leading Order}

Within the shift-symmetric Horndeski truncation, the beta functions for the
dimensionless braiding coupling $h$ and the quartic coupling $g$ can be
obtained from the $P^{-1}F$ expansion of the Wetterich equation, evaluated
on a maximally symmetric background with the York decomposition of metric
fluctuations and the Litim regulator \cite{EichhornLinoWagner2023}. To leading
order in the gravitational couplings, the results are
\begin{align}
\beta_h &= \left[3 + \frac{g}{4\pi^2}
+ G\!\left(\frac{7}{4\pi(3-4\Lambda)}
+ \frac{9}{2\pi(3-4\Lambda)^2}
- \frac{5}{9\pi(1-2\Lambda)}
- \frac{20}{9\pi(1-2\Lambda)^2}\right)
\right]h,
\label{eq:beta_h_explicit}
\\[6pt]
\beta_g &= 4g + \frac{9}{32\pi^2}g^2
+ \frac{Gg}{6\pi}\left[
\frac{8}{3-4\Lambda}
+ \frac{15}{(3-4\Lambda)^2}
- \frac{20}{(1-2\Lambda)^2}
\right] + \mathcal{O}(G^2).
\label{eq:beta_g_explicit}
\end{align}

\subsection{Structure and Interpretation of $\beta_h$}

The beta function (\ref{eq:beta_h_explicit}) has a definitive structure: $h$
appears as an overall factor on the right-hand side, confirming the symmetry
argument that $h_* = 0$ is an exact fixed point. The bracket contains three
types of contribution. The term $+3$ is the canonical scaling, arising from
the mass dimension $[h] = -3$ and the rescaling $h = k^{-3}\tilde h$. The
term $g/(4\pi^2)$ is the matter self-interaction correction from the $\chi^2$
vertex appearing in the loop diagram at order $g$. The remaining gravitational
terms arise from the transverse-traceless graviton loops (proportional to the
$(3-4\Lambda)$ denominators) and the trace scalar graviton loops (proportional
to the $(1-2\Lambda)$ denominators).

The sign of the gravitational bracket is crucial. Expanding near $G = 0$ and
$\Lambda = 0$, one finds that the gravitational contribution is negative,
meaning that gravity reduces the canonical suppression of $h$ but cannot change
the sign of the bracket without exceeding the weak-gravity bound
$G \lesssim 1.4$ \cite{EichhornLinoWagner2023,EichhornSchiffer2022handbook}.
The critical exponent $\theta_h$ at the SGFP is given by minus the value of
the bracket at $g = g_*$, $G = G_*$, $\Lambda = \Lambda_*$:
\begin{equation}
\theta_h = -\beta_h/h\big|_{\rm SGFP}
= -\left[3 + \frac{g_*}{4\pi^2} + (\text{gravitational terms})\right].
\end{equation}
Since the gravitational terms cannot make the bracket negative without
violating the weak-gravity bound, $\theta_h < 0$ throughout the accessible
region of the $(G,\Lambda)$ plane. The braiding coupling is therefore
irrelevant at all physically viable gravitational fixed points.

\subsection{Structure and Interpretation of $\beta_g$}

The beta function (\ref{eq:beta_g_explicit}) has a qualitatively different
structure. The canonical scaling term is $+4g$, reflecting the mass dimension
$[g] = -4$. The matter sector correction $9g^2/(32\pi^2)$ is the standard
quartic self-interaction. The gravitational contribution at order $G$ does not
vanish at $g = 0$: it is proportional to $Gg$, meaning that it shifts the
fixed-point value from zero but does not introduce a $g$-independent term that
would produce a non-zero fixed point from the Gaussian starting value $g = 0$
through a pure gravitational mechanism.

The fixed-point condition $\beta_g = 0$ gives
\begin{equation}
g_* = -\frac{4 + f_G(G_*,\Lambda_*)}{9/(32\pi^2)},
\end{equation}
where $f_G$ denotes the gravitational bracket in equation (\ref{eq:beta_g_explicit})
evaluated at the fixed point. For the gravitational couplings at the NGFP,
$G_* \approx 0.7$--$1.2$ and $\Lambda_* \approx 0.1$--$0.2$ (truncation
dependent), one finds $g_* < 0$ at the SGFP. The quartic coupling is therefore
shifted to a small negative value by gravity, which is
not the large positive value $g_{\rm required} \approx 9 \times 10^7\,\text{eV}^{-4}$
required for dark energy \cite{EichhornLinoWagner2023}.

%%==========================================================================
\section{Fixed-Point Analysis and the Stability Matrix}
\label{sec:stability_matrix}
%%==========================================================================

\subsection{The Stability Matrix at the SGFP}

The stability properties of the shifted Gaussian fixed point are encoded in the
stability matrix $M_{ij} = \partial\beta_{g_i}/\partial g_j|_{g_*}$ evaluated
at the fixed-point values $(h_*, g_*, G_*, \Lambda_*)$. Since $h_* = 0$ and
the beta functions factorize as $\beta_h = h\cdot f(g,G,\Lambda)$, the
$h$-direction decouples from the $(g,G,\Lambda)$ sector at the linearized
level. The stability matrix therefore has a block-diagonal structure:
\begin{equation}
M = \begin{pmatrix}
M_h & 0 \\
0 & M_{g,G,\Lambda}
\end{pmatrix},
\end{equation}
where $M_h = \partial\beta_h/\partial h|_{h_*=0} = f(g_*,G_*,\Lambda_*)$ is
a single scalar entry and $M_{g,G,\Lambda}$ is a $3\times 3$ matrix for the
remaining couplings.

\subsection{Critical Exponents}

The critical exponents are the negatives of the eigenvalues of $M$:
$\theta_I = -{\rm eig}(M)$. The critical exponent associated with $h$ is
\begin{equation}
\theta_h = -M_h = -\left[3 + \frac{g_*}{4\pi^2} + (\text{grav. bracket})\right].
\end{equation}
For representative values $G_* = 1.0$, $\Lambda_* = 0.1$, $g_* \approx -0.02$
(SGFP), the gravitational bracket in equation (\ref{eq:beta_h_explicit})
evaluates to approximately $-0.8$, giving
\begin{equation}
\theta_h \approx -(3 - 0.0005 - 0.8) \approx -2.2.
\end{equation}
This negative value confirms the irrelevance of $h$ at the SGFP. The canonical
scaling $+3$ dominates over the gravitational correction $-0.8$, leaving a net
negative critical exponent of order two \cite{EichhornLinoWagner2023}.

The RG trajectory emerging from the SGFP therefore has $h(k) \propto (k/M_{\rm Pl})^{|\theta_h|} \approx (k/M_{\rm Pl})^{2.2}$ as the scale $k$ decreases from the Planck scale. At the cosmological scale $k \sim H_0$,
\begin{equation}
h(H_0) \sim \left(\frac{H_0}{M_{\rm Pl}}\right)^{2.2}
\approx \left(10^{-61}\right)^{2.2}
\approx 10^{-134}.
\end{equation}
This is a fantastically small number. The observationally required value is
$h \sim H_0^{-3}M_{\rm Pl}^{-0} \approx 10^{126}\,{\rm eV}^{-3}$, implying a
mismatch of roughly 260 orders of magnitude in the dimensionful coupling
$\bar h = h k^3$ evaluated at $k = H_0$ \cite{EichhornLinoWagner2023}.

\subsection{Near-Critical Regime and the Weak-Gravity Bound}

Near the weak-gravity bound $G \to G_{\rm crit} \approx 1.4$ (at $\Lambda = 0$),
the gravitational bracket in equation (\ref{eq:beta_h_explicit}) increases in
magnitude. One can ask whether the bracket can become large enough to produce
$\theta_h > 0$, which would make $h$ relevant and potentially allow non-zero
braiding in the IR. A systematic scan of the accessible $(G,\Lambda)$ region
shows that this does not occur: the bracket reaches at most approximately
$-1.5$ before the matter fixed points disappear at the weak-gravity bound,
giving $\theta_h \approx -(3 - 1.5) = -1.5$ rather than zero
\cite{EichhornLinoWagner2023,EichhornSchiffer2022handbook}. The braiding
coupling remains irrelevant throughout the accessible parameter space.

%%==========================================================================
\section{The EFT-AS Gap and the Cosmological Matching Problem}
\label{sec:eft_as_gap}
%%==========================================================================

\subsection{Effective Field Theory Requirements for Dark Energy}

From the effective field theory perspective, the cosmological constant problem
requires that the quartic coupling $g$ and braiding coupling $h$ take specific
infrared values to reproduce the observed dark energy density
$\rho_\Lambda \approx (2.3\times 10^{-3}\,{\rm eV})^4 \approx 2.8\times 10^{-47}\,{\rm GeV}^4$.
In the shift-symmetric Horndeski model with $h \neq 0$, the scalar field energy
density on a de Sitter background evaluates to contributions from both $g$ and
$h$, requiring \cite{EichhornLinoWagner2023}
\begin{align}
g_{\rm EFT}
&\approx \frac{\rho_\Lambda}{H_0^4}
\approx \frac{10^{-47}\,{\rm GeV}^4}{(2.3\times 10^{-42}\,{\rm GeV})^4}
\approx 9\times 10^7\,{\rm eV}^{-4},
\\[4pt]
h_{\rm EFT}
&\approx \frac{1}{H_0^3\,M_{\rm Pl}^0}
\approx \frac{1}{(2.3\times 10^{-42}\,{\rm GeV})^3}
\approx 2.3\times 10^{33}\,{\rm eV}^{-3}.
\end{align}
These are enormous numbers in natural units. They correspond to the intrinsic
scale of the cosmological constant fine-tuning problem, expressed now as
specific requirements on Horndeski coupling values.

\subsection{Asymptotic-Safety Predictions versus EFT Requirements}

The asymptotic-safety prediction for the RG-evolved coupling $g$ at the
cosmological scale is obtained by running the quartic coupling from the
Planck scale $k = M_{\rm Pl}$ to the Hubble scale $k = H_0$. Using the
linearized solution near the SGFP:
\begin{equation}
g(H_0)
\approx g_* + C_g\left(\frac{H_0}{M_{\rm Pl}}\right)^{|\theta_g|}
\approx g_* + C_g\times 10^{-244}.
\end{equation}
The integration constant $C_g$ is of order one for trajectories on the
UV-critical surface. The running therefore produces
$g(H_0) \approx g_* \approx -0.02$, in natural units where the RG scale
sets the mass scale, which translates to the dimensionful coupling
$\bar g(H_0) = g_* H_0^4 \approx -0.02\times (2.3\times 10^{-42}\,{\rm GeV})^4
\approx -6\times 10^{-49}\,{\rm GeV}^4$.

Converting to the observational units used in the EFT comparison:
\begin{equation}
g_{\rm AS}(H_0) \approx g_* \times H_0^{-4}
\approx -0.02\times (10^{-42}\,{\rm GeV})^{-4}
\approx -2\times 10^{79}\,{\rm GeV}^{-4}
\approx -2\times 10^{-7}\,{\rm eV}^{-4}.
\end{equation}
Comparing with $g_{\rm EFT} \approx 9\times 10^7\,{\rm eV}^{-4}$, the
mismatch is a factor of $\sim 10^{15}$. The sign is also wrong: asymptotic
safety predicts a small negative $g_*$, while the EFT requirement is a large
positive $g_{\rm EFT}$ \cite{EichhornLinoWagner2023}.

For the braiding coupling the situation is even starker. Asymptotic safety
predicts $h(H_0) \approx 0$ to all practical purposes, while the EFT
requirement is $h_{\rm EFT} \approx 2.3\times 10^{33}\,{\rm eV}^{-3}$.

\subsection{Implications for the Hubble Tension Resolution}

The EFT-AS gap has direct implications for proposed resolutions of the Hubble
tension. Models in which the Horndeski braiding parameter $\alpha_B \neq 0$
generates a time-varying dark energy equation of state can shift the CMB-inferred
$H_0$ toward the locally measured value. However, such models require
$h_{\rm EFT} \sim 10^{33}\,{\rm eV}^{-3}$, whereas asymptotic safety predicts
$h(H_0) \sim 10^{-134}$ in the same units---a discrepancy of 167 orders of
magnitude. No modification to the polynomial truncation or to the regulator
scheme is capable of bridging a gap of this magnitude
\cite{EichhornLinoWagner2023,EichhornHebeckerPawlowskiWalcher2024}. The
Hubble tension cannot be resolved through shift-symmetric kinetic braiding
in a UV-complete theory.

The $\sigma_8$ tension, which involves the growth rate of structure, faces the
same obstruction: modifications to the gravitational slip parameter through
$\alpha_B \neq 0$ require the same non-zero braiding coupling that asymptotic
safety forces to zero. Both cosmological tensions therefore remain unaddressed
by the simplest asymptotically safe Horndeski models \cite{Planck2018,Riess2019}.

%%==========================================================================
\section{Gravitational Wave Signatures and the Observational Bridge}
\label{sec:gw_signatures}
%%==========================================================================

\subsection{Cosmic String Networks and the Stochastic Background}

A cosmic string network formed at a symmetry-breaking scale $v$ generates a
stochastic gravitational-wave background through the oscillation and decay of
string loops. The energy spectrum of this background, expressed as the
fractional energy density per logarithmic frequency interval, takes the form
\cite{NANOGrav2023,Maggiore2008}
\begin{equation}
\Omega_{\rm GW}(f)
= \frac{8\pi}{3H_0^2}f\,\frac{dE_{\rm GW}}{dV\,df},
\end{equation}
where the frequency-dependent emission rate depends on the string tension
$\mu = 2\pi v^2\ln(\lambda_4/g_{\rm g}^2)^{1/2}$, the loop number density, and
the loop lifetime. In the relevant frequency band of pulsar timing arrays
$f \sim 1\,{\rm yr}^{-1} \approx 3\times 10^{-8}\,{\rm Hz}$, the spectrum
of a network of Nambu-Goto strings in the scaling regime scales approximately
as a power law
\begin{equation}
\Omega_{\rm GW}(f) \propto (G\mu)^{x}\,f^{n},
\end{equation}
where $G\mu$ is the dimensionless string tension (Newton's constant times string
tension), $x \approx 1$, and $n$ depends on the loop distribution model
\cite{NANOGrav2023}.

The NANOGrav 15-year dataset reports a detection of the stochastic background
with amplitude and spectral index \cite{NANOGrav2023}
\begin{equation}
h_c(f) \approx A_{\rm GW}\left(\frac{f}{f_{\rm yr}}\right)^{(3-\gamma)/2},
\quad A_{\rm GW} \approx 2.4\times 10^{-15},
\quad \gamma \approx 3.2,
\end{equation}
where $f_{\rm yr} = 1\,{\rm yr}^{-1}$. The Hellings-Downs angular correlation
function, characteristic of a gravitational-wave background, was detected at
$\sim 4\sigma$ significance. However, the spectral shape and amplitude are
consistent with a population of supermassive black-hole binaries and
inconsistent with the spectrum expected from a stable Nambu-Goto string
network at the $v \lesssim 4.6\times 10^{14}\,{\rm GeV}$ scale permitted by
observational bounds \cite{NANOGrav2023,EichhornLinoMiqueleto2023}.

\subsection{Connecting String Tension to the Symmetry-Breaking Scale}

The constraint from NANOGrav on the symmetry-breaking scale can be derived
by requiring that the theoretical prediction for $\Omega_{\rm GW}$ not exceed
the observed signal at the relevant frequencies. For a stable cosmic string
network the constraint translates to an upper bound on the string tension
\begin{equation}
G\mu \lesssim 2\times 10^{-10},
\end{equation}
which corresponds to a symmetry-breaking scale
\begin{equation}
v \lesssim 4.6\times 10^{14}\,{\rm GeV} \approx 3.8\times 10^{-5}M_{\rm Pl}.
\label{eq:v_bound}
\end{equation}
The asymptotic-safety prediction for the minimal Abelian Higgs model is
$v \sim M_{\rm Pl} = 1.22\times 10^{19}\,{\rm GeV}$ \cite{EichhornLinoMiqueleto2023}.
The ratio of the predicted and allowed scales is
\begin{equation}
\frac{v_{\rm AS}}{v_{\rm max}}
\approx \frac{M_{\rm Pl}}{3.8\times 10^{-5}M_{\rm Pl}}
\approx 2.6\times 10^4.
\end{equation}
This is the scale mismatch that the forty-two fermion extension must address:
the fermions must screen the gauge coupling enough to shift $v$ downward by
over four orders of magnitude from the Planck scale.

\subsection{The AS-GW Interface as a Precision Test}

The chain of reasoning connecting asymptotic safety to the gravitational
wave background is a closed logical loop: UV fixed-point structure constrains
$v$ through the gauge coupling fixed point, $v$ determines $G\mu$, $G\mu$
determines the amplitude and spectral shape of $\Omega_{\rm GW}(f)$, and
$\Omega_{\rm GW}(f)$ is constrained by NANOGrav. As the sensitivity of pulsar
timing arrays improves---with the Square Kilometre Array (SKA) expected to
reach $G\mu \sim 10^{-12}$---the constraint (\ref{eq:v_bound}) will tighten,
further probing the parameter space of UV-complete cosmic string models
\cite{EichhornLinoMiqueleto2023,BuoninfanteEtal2025}. Either a signal
consistent with a sub-Planckian string scale is detected, requiring the
forty-two fermion extension to be part of the true theory, or no string
signal is found, consistent with the minimal AS prediction that cosmic strings
are a Planck-scale phenomenon invisible to current and near-future PTAs.

\subsection{Tensor Modes and the Speed of Gravitational Waves}

The GW170817 constraint on the tensor propagation speed, $|c_t - 1| \lesssim
10^{-15}$ \cite{Abbott2016}, provides a separate and complementary constraint
on the Horndeski sector. Within the Horndeski framework the tensor speed is
\cite{Clifton2012,Nicolis2009}
\begin{equation}
c_t^2 = 1 + \frac{2XG_{4,X}}{M_{\rm Pl}^2 + 2XG_{4,X}},
\end{equation}
where $G_{4,X} = \partial G_4/\partial X$ is the derivative of the $L_4$
Horndeski function with respect to the kinetic variable $X = -\chi$.
The GW170817 constraint requires $G_{4,X} \approx 0$, which eliminates the
$L_4$ and $L_5$ sectors and restricts viable Horndeski theories to the
$L_2$-$L_3$ subspace \cite{EichhornLinoWagner2023}. This is precisely the
shift-symmetric kinetic braiding sector analyzed in this work. The
constraint $c_t = 1$ is therefore not an independent result but a
consistency condition that uniquely selects the sector in which the
asymptotic-safety analysis applies.

%%==========================================================================
\section{Conclusion}
%%==========================================================================

The asymptotic-safety paradigm, implemented through the Functional
Renormalization Group, provides a powerful top-down filter on the space of
phenomenological models for physics beyond the Standard Model. We have reviewed
in detail how this filter operates for two model classes of direct observational
relevance: shift-symmetric Horndeski theories with kinetic braiding, and the
Abelian Higgs model as a source of cosmic strings.

For kinetic braiding, the conclusion is definitive within the polynomial
truncation: the braiding coupling $h$ is symmetry-protected and vanishes
identically at the shifted Gaussian fixed point. Its IR value is predicted to
be zero, which is incompatible with the non-zero braiding required for dynamical
dark energy \cite{EichhornLinoWagner2023}. The quartic coupling $g$ is
similarly suppressed by approximately fifty-seven orders of magnitude relative
to the value needed to match the observed dark energy density. Kinetic braiding
models are swampland candidates: they cannot be connected to an asymptotically
safe UV completion without abandoning the phenomenological features that
motivate them \cite{EichhornHebeckerPawlowskiWalcher2024}.

For cosmic strings in the Abelian Higgs model, the conclusion is more
conditional: the minimal model fails, but an extended model with at least
forty-two charged Dirac fermions may satisfy the UV consistency conditions
\cite{EichhornLinoMiqueleto2023}. This forty-two fermion barrier quantifies
the complexity tax imposed by asymptotic safety on the cosmic string hypothesis.
The recent NANOGrav results \cite{NANOGrav2023}, which disfavor a stable-string
interpretation of the observed stochastic gravitational-wave background in
favor of supermassive black-hole binaries, provide independent observational
support for the theoretical exclusion of the minimal model.

These results collectively support a picture of the asymptotically safe
landscape as a small landscape: the UV consistency requirement reduces the
apparently vast space of possible low-energy effective theories to a much
smaller set of genuinely viable candidates \cite{Eichhorn2019,Percacci2017,
BuoninfanteEtal2025}. This reduction is not a limitation of the framework but
its defining virtue. A theory that cannot explain the origin of its own
parameters from first principles is not a fundamental theory; it is an
approximation \cite{WilsonKogut1974,Weinberg1979}. The asymptotic-safety
program, by demanding that every coupling either flow to a fixed-point value
or be genuinely measurable, transforms the question of model selection from a
matter of phenomenological preference into a matter of mathematical necessity.

The program is not complete. Non-polynomial truncations, Lorentzian extensions,
and more general geometric frameworks remain open directions that may modify
or extend the conclusions reviewed here \cite{BorissovaDittrichEichhornSchiffer2025,
EichhornGurauGyftopoulos2025}. But the methodology is established, the tools
are in place, and the first generation of results is already shaping the
landscape of viable beyond-Standard-Model physics in ways that decades of
bottom-up phenomenology alone could not achieve \cite{EichhornSchiffer2022handbook,
Eichhorn2022status}. The universe we observe may ultimately be far more
constrained---and far more elegant---than the space of effective theories would
suggest.

%%==========================================================================
%% APPENDICES
%%==========================================================================

\newpage
\section*{Appendices}

\appendix

%%==========================================================================
\section{Derivation of the Wetterich Flow Equation}
\label{app:wetterich}
%%==========================================================================

The central object of the Functional Renormalization Group is the
scale-dependent effective action $\Gamma_k$. Its flow with respect to the
infrared cutoff scale $k$ is governed by the Wetterich equation \cite{Wetterich1993}.

We begin with the partition function in the presence of a regulator term that
suppresses infrared modes:
\begin{equation}
Z_k[J] =
\int \mathcal{D}\Phi \;\exp\!\left(
- S[\Phi] - \Delta S_k[\Phi] + J\cdot\Phi
\right),
\end{equation}
where $\Phi$ collectively denotes all fields and the regulator term is
$\Delta S_k = \tfrac{1}{2}\Phi R_k \Phi$. The generating functional of
connected correlators is $W_k[J] = \ln Z_k[J]$, and the classical field is
defined by $\Phi_{\rm cl} = \delta W_k/\delta J$. The scale-dependent effective
action is obtained through a modified Legendre transform:
\begin{equation}
\Gamma_k[\Phi_{\rm cl}]
= J\cdot\Phi_{\rm cl} - W_k[J] - \tfrac{1}{2}\Phi_{\rm cl} R_k \Phi_{\rm cl}.
\end{equation}
From the definition of $\Gamma_k$ and the stationarity condition
$\delta\Gamma_k/\delta\Phi_{\rm cl} = J - R_k\Phi_{\rm cl}$, one obtains
the second functional derivative relation
\begin{equation}
\left(\Gamma_k^{(2)} + R_k\right)_{xy}
= \left(W_k^{(2)}\right)_{xy}^{-1}.
\end{equation}
Differentiating the modified Legendre transform with respect to the scale
parameter $t = \ln k$ and using the chain rule on both sides, one arrives at
the exact Wetterich flow equation \cite{Wetterich1993,Berges2002}:
\begin{equation}
\partial_t \Gamma_k
=
\frac{1}{2}\,
\mathrm{Tr}
\left[
(\Gamma_k^{(2)} + R_k)^{-1}
\partial_t R_k
\right].
\end{equation}
This equation has the structural form of a one-loop expression but is exact:
all quantum corrections are encoded in the full propagator $(\Gamma_k^{(2)}
+ R_k)^{-1}$, which itself depends on $\Gamma_k$ \cite{Dupuis2021,Pawlowski2007}.

%%==========================================================================
\section{York Decomposition of Metric Fluctuations}
\label{app:york}
%%==========================================================================

In functional renormalization group calculations for gravity, it is convenient
to decompose metric fluctuations into irreducible components under
diffeomorphisms \cite{Reuter1998,Percacci2017}.

The background-field method writes the full metric as
$g_{\mu\nu} = \bar g_{\mu\nu} + h_{\mu\nu}$. The York decomposition expresses
$h_{\mu\nu}$ as
\begin{equation}
h_{\mu\nu}
=
h_{\mu\nu}^{\rm TT}
+\nabla_\mu \xi_\nu + \nabla_\nu \xi_\mu
+\nabla_\mu \nabla_\nu \sigma
-\tfrac{1}{4}\bar g_{\mu\nu} \nabla^2 \sigma
+\tfrac{1}{4}\bar g_{\mu\nu} h,
\end{equation}
where the transverse-traceless (TT) modes satisfy
\begin{equation}
\nabla^\mu h_{\mu\nu}^{\rm TT} = 0, \qquad
\bar g^{\mu\nu} h_{\mu\nu}^{\rm TT} = 0.
\end{equation}
The decomposition therefore separates the metric fluctuations into spin-2
transverse-traceless modes $h_{\mu\nu}^{\rm TT}$ carrying the propagating
graviton degrees of freedom, transverse vector modes $\xi_\mu$, and two
scalar modes $\sigma$ and $h$. In the Landau-DeWitt gauge the vector and
longitudinal scalar modes decouple from the physical spectrum. The TT sector
typically dominates the RG flow of the gravitational couplings, while the
trace scalar mode couples to matter fields through non-minimal interaction
operators such as $\zeta\chi R$ \cite{EichhornLinoWagner2023}.

%%==========================================================================
\section{Heat-Kernel Evaluation of Functional Traces}
\label{app:heatkernel}
%%==========================================================================

Functional traces appearing in the Wetterich equation can be evaluated using
heat-kernel techniques \cite{Berges2002,Dupuis2021}. Consider an operator
of Laplace type $\Delta = -\nabla^2 + U$. The trace of a function of this
operator can be written as
\begin{equation}
\mathrm{Tr}\, f(\Delta)
=
\int d^4x \sqrt{g}
\sum_{n=0}^{\infty}
Q_{4-n}[f]\, a_n(\Delta),
\end{equation}
where the $a_n$ are the Seeley-DeWitt heat-kernel coefficients, and the
$Q$-functionals are defined by
\begin{equation}
Q_n[f]
=
\frac{1}{\Gamma(n)}
\int_0^\infty ds\, s^{n-1} f(s).
\end{equation}
The leading coefficients are $a_0 = 1/(4\pi)^2$, $a_2 = R/(6(4\pi)^2)$,
and higher terms involving curvature invariants. For the optimized Litim
regulator $R_k(p^2) = (k^2 - p^2)\Theta(k^2 - p^2)$, the $Q$-functionals
reduce to simple rational expressions:
\begin{equation}
Q_n\!\left[\frac{\partial_t R_k}{\Gamma_k^{(2)}+R_k}\right]
=
\frac{2k^{2n}}{(k^2 + m^2)^p},
\end{equation}
for appropriate powers $p$ depending on $n$ and the regulator form. Using
these identities, the functional trace in the Wetterich equation can be
expressed as a polynomial in the curvature invariants and field-dependent
masses, allowing the beta functions of gravitational and matter couplings to
be extracted systematically \cite{Pawlowski2007,Berges2002}.

%%==========================================================================
\section{Projection of Beta Functions}
\label{app:projection}
%%==========================================================================

To determine the renormalization group flow of specific couplings, one
projects the flow equation onto the corresponding operators \cite{Pawlowski2007,
Dupuis2021}. Consider a truncated effective action
\begin{equation}
\Gamma_k = \sum_i g_i(k) \int d^4x \sqrt{g}\, \mathcal{O}_i.
\end{equation}
Applying the flow equation yields $\partial_t \Gamma_k =
\tfrac{1}{2}\,\mathrm{Tr}[(\Gamma_k^{(2)} + R_k)^{-1}\partial_t R_k]$.
To obtain the beta function for a coupling $g_i$, one isolates the coefficient
of the operator $\mathcal{O}_i$:
\begin{equation}
\beta_{g_i} = \frac{\partial}{\partial \mathcal{O}_i}\left(\partial_t \Gamma_k\right).
\end{equation}
In practice this projection is implemented by evaluating the flow on specific
background field configurations. For example, to extract the flow of the
Einstein-Hilbert term, one evaluates the trace on a maximally symmetric
background where
\begin{equation}
R_{\mu\nu\rho\sigma}
= \frac{R}{12}(g_{\mu\rho}g_{\nu\sigma} - g_{\mu\sigma}g_{\nu\rho}).
\end{equation}
Scalar-sector couplings are obtained by expanding around constant
scalar-field backgrounds and identifying the coefficients of the corresponding
operators. For the Horndeski system, the beta functions $\beta_h$ and $\beta_g$
are extracted by projecting onto terms proportional to $\chi\Box\phi$ and
$\chi^2$ respectively, using the $P^{-1}F$ expansion described in
Section~\ref{sec:method} \cite{EichhornLinoWagner2023}.

%%==========================================================================
\section{Renormalization Group Trajectories Near the Fixed Point}
\label{app:stability}
\label{app:rgtrajectories}
%%==========================================================================

Close to the ultraviolet fixed point, the RG flow is governed by the
linearized equation (see Section~3.1)
\begin{equation}
k\partial_k \delta g_i = B_{ij}\,\delta g_j,
\end{equation}
where $B_{ij} = \partial\beta_i/\partial g_j|_{g_*}$ is the stability matrix
and $\delta g_i = g_i - g_i^*$ denotes the deviation from the fixed-point
values. Diagonalizing the matrix $B$ yields eigenvectors $v_I$ and eigenvalues
$-\theta_I$. The general solution is
\begin{equation}
g_i(k)
= g_i^*
+ \sum_I C_I v_I^{(i)}
\left(\frac{k}{k_0}\right)^{-\theta_I},
\end{equation}
where $k_0$ is a reference scale and $C_I$ are integration constants
\cite{WilsonKogut1974,ReuterSaueressig2012}. The constants $C_I$ associated
with relevant directions ($\theta_I > 0$) are the free parameters of the
theory, to be fixed by experiment; those associated with irrelevant directions
($\theta_I < 0$) are driven to zero as $k \to \infty$ and are therefore
predicted by the UV fixed point.

For the shifted Gaussian fixed point of the Horndeski system, the critical
exponent $\theta_h$ associated with the braiding coupling satisfies
$\theta_h < 0$, confirming the irrelevance of $h$. As $k$ decreases from the
Planck scale toward cosmological scales, $h(k)$ remains exponentially
suppressed:
\begin{equation}
h(k) \sim h_* + C_h \left(\frac{k}{M_{\rm Pl}}\right)^{|\theta_h|},
\end{equation}
with $h_* = 0$. The infrared value $h(H_0)$ is therefore suppressed by the
ratio $(H_0/M_{\rm Pl})^{|\theta_h|} \approx 10^{-60|\theta_h|}$, making it
completely negligible for cosmological phenomenology \cite{EichhornLinoWagner2023}.

%%==========================================================================
\section{Operator Dimension Table and Truncation Systematics}
\label{app:dimensions}
%%==========================================================================

The operators entering the FRG truncation can be organized systematically
according to their canonical mass dimension $[{\mathcal O}]$ in four Euclidean
dimensions \cite{WilsonKogut1974,Peskin1995}. For a $d$-dimensional spacetime
integral $\int d^4x\,\sqrt{g}\,\mathcal{O}$, the action has dimension zero,
so the coupling constant $g_{\mathcal O}$ has dimension $4 - [\mathcal{O}]$.
Operators with $[\mathcal{O}] > 4$ are power-counting irrelevant.

For the scalar field in four dimensions, the canonical dimension is
$[\phi] = 1$, giving $[\chi] = [\partial\phi\,\partial\phi] = 4$. The coupling
to the standard kinetic term $\int\sqrt{g}\,\chi$ is therefore dimensionless.
The cubic operator $\chi\Box\phi$ has dimension $[\chi\Box\phi] = 4 + 1 + 2 = 7$,
giving $[h] = 4 - 7 = -3$. The quartic operator $\chi^2$ has dimension
$[\chi^2] = 8$, giving $[g] = 4 - 8 = -4$. These values confirm the
power-counting irrelevance of both operators in the absence of gravity.

The gravitational sector introduces the Newton coupling $G_N = M_{\rm Pl}^{-2}$
with dimension $[G_N] = -2$ in four dimensions, or equivalently the
dimensionless combination $g = G_N k^2$ used throughout this essay. The
cosmological constant $\bar\Lambda$ has dimension $[\bar\Lambda] = 2$,
corresponding to the dimensionless $\Lambda = \bar\Lambda k^{-2}$.

The beta functions in terms of dimensionless couplings are obtained by
substituting $g_{\mathcal O} = k^{4-[\mathcal O]}\,\tilde g_{\mathcal O}$
and differentiating with respect to $t = \ln k$, yielding equation
(\ref{eq:dimless_rg}) with the canonical scaling term $-d_i\tilde g_i$
where $d_i = 4 - [\mathcal{O}_i]$. For $h$ and $g$, this gives canonical
scaling exponents of $+3$ and $+4$ respectively---the ``restoring force''
that drives these couplings toward zero in the absence of quantum corrections.
The gravitational antiscreening contribution $\beta_g^{\rm grav}$ must
overcome this canonical scaling by an amount sufficient to produce $g_* \neq 0$,
which it does for any $G > 0$, while for $h$ the symmetry argument ensures
that $\beta_h^{\rm grav} \propto h$ and cannot override the symmetry protection
\cite{EichhornLinoWagner2023,EichhornSchiffer2022handbook}.

\newpage

\begin{thebibliography}{120}
\sloppy

%% --- Renormalization and Quantum Field Theory ---

\bibitem{Wilson1971}
K.~G.~Wilson,
\textit{Renormalization Group and Critical Phenomena},
Phys.\ Rev.\ B \textbf{4}, 3174 (1971).

\bibitem{WilsonKogut1974}
K.~G.~Wilson and J.~Kogut,
\textit{The Renormalization Group and the $\epsilon$ Expansion},
Phys.\ Rep.\ \textbf{12}, 75--200 (1974).

\bibitem{Polchinski1984}
J.~Polchinski,
\textit{Renormalization and Effective Lagrangians},
Nucl.\ Phys.\ B \textbf{231}, 269--295 (1984).

\bibitem{Wegner1972}
F.~J.~Wegner,
\textit{Corrections to Scaling Laws},
Phys.\ Rev.\ B \textbf{5}, 4529 (1972).

\bibitem{ZinnJustin2002}
J.~Zinn-Justin,
\textit{Quantum Field Theory and Critical Phenomena},
Oxford University Press, Oxford (2002).

\bibitem{Cardy1996}
J.~Cardy,
\textit{Scaling and Renormalization in Statistical Physics},
Cambridge University Press, Cambridge (1996).

\bibitem{Peskin1995}
M.~Peskin and D.~Schroeder,
\textit{An Introduction to Quantum Field Theory},
Addison-Wesley, Reading (1995).

\bibitem{Weinberg1995}
S.~Weinberg,
\textit{The Quantum Theory of Fields}, Vols.\ I--III,
Cambridge University Press, Cambridge (1995).

%% --- Functional Renormalization Group ---

\bibitem{Wetterich1993}
C.~Wetterich,
\textit{Exact Evolution Equation for the Effective Potential},
Phys.\ Lett.\ B \textbf{301}, 90--94 (1993).

\bibitem{Berges2002}
J.~Berges, N.~Tetradis, and C.~Wetterich,
\textit{Non-Perturbative Renormalization Flow in Quantum Field Theory
and Statistical Physics},
Phys.\ Rep.\ \textbf{363}, 223--386 (2002).

\bibitem{Pawlowski2007}
J.~M.~Pawlowski,
\textit{Aspects of the Functional Renormalisation Group},
Annals Phys.\ \textbf{322}, 2831 (2007).

\bibitem{Dupuis2021}
N.~Dupuis et al.,
\textit{The Nonperturbative Functional Renormalization Group},
Phys.\ Rep.\ \textbf{910}, 1--114 (2021).

%% --- Asymptotic Safety and Quantum Gravity ---

\bibitem{Weinberg1979}
S.~Weinberg,
\textit{Ultraviolet Divergences in Quantum Theories of Gravitation},
in S.~W.~Hawking and W.~Israel (eds.),
\textit{General Relativity: An Einstein Centenary Survey},
Cambridge University Press, Cambridge (1979).

\bibitem{Reuter1998}
M.~Reuter,
\textit{Nonperturbative Evolution Equation for Quantum Gravity},
Phys.\ Rev.\ D \textbf{57}, 971--985 (1998).

\bibitem{ReuterSaueressig2012}
M.~Reuter and F.~Saueressig,
\textit{Quantum Einstein Gravity},
New J.\ Phys.\ \textbf{14}, 055022 (2012).

\bibitem{Percacci2017}
R.~Percacci,
\textit{An Introduction to Covariant Quantum Gravity and Asymptotic Safety},
World Scientific, Singapore (2017).

\bibitem{Eichhorn2019}
A.~Eichhorn,
\textit{An Asymptotically Safe Guide to Quantum Gravity and Matter},
Front.\ Astron.\ Space Sci.\ \textbf{5}, 47 (2019).

\bibitem{EichhornSchiffer2022handbook}
A.~Eichhorn and M.~Schiffer,
\textit{Asymptotic Safety of Gravity with Matter},
in \textit{Handbook of Quantum Gravity},
Springer Nature Singapore (2022).

\bibitem{Bonanno2002}
A.~Bonanno and M.~Reuter,
\textit{Cosmology with Self-Adjusting Vacuum Energy Density from a
Renormalization Group Fixed Point},
Phys.\ Lett.\ B \textbf{527}, 9 (2002).

%% --- Statistical Physics and Coarse-Graining ---

\bibitem{Goldenfeld1992}
N.~Goldenfeld,
\textit{Lectures on Phase Transitions and the Renormalization Group},
Addison-Wesley, Reading (1992).

\bibitem{Kadanoff1966}
L.~P.~Kadanoff,
\textit{Scaling Laws for Ising Models Near $T_c$},
Physics \textbf{2}, 263 (1966).

\bibitem{Stanley1971}
H.~Stanley,
\textit{Introduction to Phase Transitions and Critical Phenomena},
Oxford University Press, Oxford (1971).

%% --- Information Geometry ---

\bibitem{Fisher1922}
R.~A.~Fisher,
\textit{On the Mathematical Foundations of Theoretical Statistics},
Phil.\ Trans.\ R.\ Soc.\ A \textbf{222}, 309--368 (1922).

\bibitem{AmariNagaoka2000}
S.~Amari and H.~Nagaoka,
\textit{Methods of Information Geometry},
American Mathematical Society, Providence (2000).

\bibitem{Amari2016}
S.~Amari,
\textit{Information Geometry and Its Applications},
Springer, Berlin (2016).

\bibitem{Ruppeiner1995}
G.~Ruppeiner,
\textit{Riemannian Geometry in Thermodynamic Fluctuation Theory},
Rev.\ Mod.\ Phys.\ \textbf{67}, 605 (1995).

%% --- Metric Geometry ---

\bibitem{Gromov1981}
M.~Gromov,
\textit{Groups of Polynomial Growth and Expanding Maps},
Publ.\ Math.\ IH\'ES \textbf{53}, 53--73 (1981).

\bibitem{Burago2001}
D.~Burago, Y.~Burago, and S.~Ivanov,
\textit{A Course in Metric Geometry},
American Mathematical Society, Providence (2001).

%% --- Category Theory and Higher Structures ---

\bibitem{MacLane1998}
S.~Mac~Lane,
\textit{Categories for the Working Mathematician},
Springer, New York (1998).

\bibitem{Awodey2010}
S.~Awodey,
\textit{Category Theory},
Oxford University Press, Oxford (2010).

\bibitem{Lurie2009}
J.~Lurie,
\textit{Higher Topos Theory},
Princeton University Press, Princeton (2009).

\bibitem{Lurie2017}
J.~Lurie,
\textit{Higher Algebra},
available at \url{https://www.math.ias.edu/~lurie/} (2017).

%% --- Scalar-Tensor Gravity and Modified Gravity ---

\bibitem{Horndeski1974}
G.~W.~Horndeski,
\textit{Second-Order Scalar-Tensor Field Equations in a
Four-Dimensional Space},
Int.\ J.\ Theor.\ Phys.\ \textbf{10}, 363--384 (1974).

\bibitem{Nicolis2009}
A.~Nicolis, R.~Rattazzi, and E.~Trincherini,
\textit{The Galileon as a Local Modification of Gravity},
Phys.\ Rev.\ D \textbf{79}, 064036 (2009).

\bibitem{Clifton2012}
T.~Clifton et al.,
\textit{Modified Gravity and Cosmology},
Phys.\ Rep.\ \textbf{513}, 1 (2012).

%% --- Cosmology and Structure Formation ---

\bibitem{Peebles1980}
P.~J.~E.~Peebles,
\textit{The Large-Scale Structure of the Universe},
Princeton University Press, Princeton (1980).

\bibitem{Mukhanov2005}
V.~Mukhanov,
\textit{Physical Foundations of Cosmology},
Cambridge University Press, Cambridge (2005).

\bibitem{Dodelson2003}
S.~Dodelson,
\textit{Modern Cosmology},
Academic Press, San Diego (2003).

\bibitem{Weinberg2008}
S.~Weinberg,
\textit{Cosmology},
Oxford University Press, Oxford (2008).

%% --- Observational Cosmology ---

\bibitem{Planck2018}
Planck Collaboration,
\textit{Planck 2018 Results. VI. Cosmological Parameters},
Astron.\ Astrophys.\ \textbf{641}, A6 (2020).

\bibitem{Riess2019}
A.~Riess et al.,
\textit{Large Magellanic Cloud Cepheid Standards Provide a 1\%
Foundation for the Determination of the Hubble Constant},
Astrophys.\ J.\ \textbf{876}, 85 (2019).

%% --- Gravitational Waves ---

\bibitem{NANOGrav2023}
NANOGrav Collaboration,
\textit{The NANOGrav 15-year Data Set: Evidence for a
Gravitational-Wave Background},
Astrophys.\ J.\ Lett.\ \textbf{951}, L8 (2023).

\bibitem{Abbott2016}
B.~P.~Abbott et al.,
\textit{Observation of Gravitational Waves from a Binary Black Hole Merger},
Phys.\ Rev.\ Lett.\ \textbf{116}, 061102 (2016).

\bibitem{Maggiore2008}
M.~Maggiore,
\textit{Gravitational Waves: Volume 1: Theory and Experiments},
Oxford University Press, Oxford (2008).

%% --- Asymptotically Safe Gravity-Matter Systems: Research Front ---

\bibitem{Eichhorn2022status}
A.~Eichhorn,
\textit{Status Update: Asymptotically Safe Gravity-Matter Systems},
Nuovo Cim.\ C \textbf{45}, 21 (2022).

\bibitem{EichhornRay2024}
A.~Eichhorn and S.~Ray,
\textit{Suppression of Proton Decay in Quantum Gravity},
Phys.\ Lett.\ B \textbf{850}, 138529 (2024).

\bibitem{EichhornLinoMiqueleto2023}
A.~Eichhorn, R.~R.~Lino dos Santos, and J.~L.~Miqueleto,
\textit{From Quantum Gravity to Gravitational Waves through Cosmic Strings},
Phys.\ Rev.\ D \textbf{109}, 026013 (2023).

\bibitem{deBritoEichhornRay2023}
G.~P.~de~Brito, A.~Eichhorn, and S.~Ray,
\textit{Light Fermions in Color: Why the Quark Mass is Not the Planck Mass},
arXiv:2311.16066 (2023).

\bibitem{EichhornHebeckerPawlowskiWalcher2024}
A.~Eichhorn, A.~Hebecker, J.~M.~Pawlowski, and J.~Walcher,
\textit{The Absolute Swampland},
EPL \textbf{149}, 39001 (2024).

\bibitem{EichhornGurauGyftopoulos2025}
A.~Eichhorn, R.~Gurau, and Z.~Gyftopoulos,
\textit{Asymptotic Safety Meets Tensor Field Theory: Towards a New Class of
Gravity-Matter Systems},
Phys.\ Rev.\ D \textbf{111}, 085030 (2025).

\bibitem{deBritoEichhornPereiraYamada2025}
G.~P.~de~Brito, A.~Eichhorn, A.~D.~Pereira, and M.~Yamada,
\textit{Neutrino Mass Generation in Asymptotically Safe Gravity},
Phys.\ Rev.\ D \textbf{112}, 066004 (2025).

\bibitem{EichhornSchifferCPT2025}
A.~Eichhorn and M.~Schiffer,
\textit{No Dynamical CPT Symmetry Restoration in Quantum Gravity},
arXiv:2506.12001 (2025).

\bibitem{EichhornHeld2025}
A.~Eichhorn and A.~Held,
\textit{Initial Distribution at the Planck Scale for the
Third-Generation Standard Model},
arXiv:2506.12548 (2025).

\bibitem{BorissovaDittrichEichhornSchiffer2025}
J.~Borissova, B.~Dittrich, A.~Eichhorn, and M.~Schiffer,
\textit{Renormalization Group Flows in Area-Metric Gravity},
arXiv:2507.02034 (2025).

\bibitem{EichhornGyftopoulos2025}
A.~Eichhorn, Z.~Gyftopoulos, and A.~Held,
\textit{Quark and Lepton Mixing in the Asymptotically Safe Standard Model},
arXiv:2507.18304 (2025).

\bibitem{EichhornFernandes2025}
A.~Eichhorn and P.~G.~S.~Fernandes,
\textit{Regular Black Holes without Mass-Inflation Instability and
Gravastars from Modified Gravity},
arXiv:2508.00686 (2025).

\bibitem{AssantEichhornKnorr2025}
G.~Assant, A.~Eichhorn, and B.~Knorr,
\textit{Towards Theory Constraints on Ultralight Dark Matter from
Quantum Gravity},
arXiv:2510.23808 (2025).

\bibitem{DelaportEichhorn2025}
H.~Delaporte and A.~Eichhorn,
\textit{The Principled-Parameterized Approach to Gravitational Collapse},
JCAP \textbf{03}, 074 (2025).

\bibitem{EichhornSchifferPedersen2025}
A.~Eichhorn, M.~Schiffer, and A.~O.~Pedersen,
\textit{Application of Positivity Bounds in Asymptotically Safe Gravity},
Eur.\ Phys.\ J.\ C \textbf{85}, 733 (2025).

\bibitem{BuoninfanteEtal2025}
L.~Buoninfante et al.,
\textit{Visions in Quantum Gravity},
SciPost Phys.\ Comm.\ Rep.\ \textbf{11} (2025).

\bibitem{CastroEichhornGurau2026}
A.~Castro, A.~Eichhorn, and R.~Gurau,
\textit{Towards a Quantitative Characterization of Gravitational
Universality Classes for Order-4 Random Tensor Models},
arXiv:2602.09257 (2026).

\bibitem{EichhornFernandes2026}
A.~Eichhorn and P.~G.~S.~Fernandes,
\textit{Scalar-Gauss-Bonnet Gravity with a Cubic Coupling},
arXiv:2603.05064 (2026).

\bibitem{deBritoEichhornPfeiffer2023a}
G.~P.~de~Brito, A.~Eichhorn, and C.~Pfeiffer,
\textit{Towards a Bound on the Higgs Mass in Causal Set Quantum Gravity},
Gen.\ Relativ.\ Gravit.\ \textbf{55}, 128 (2023).

\bibitem{deBritoEichhornPfeiffer2023b}
G.~P.~de~Brito, A.~Eichhorn, and C.~Pfeiffer,
\textit{Higher-Order Curvature Operators in Causal Set Quantum Gravity},
Eur.\ Phys.\ J.\ Plus \textbf{138}, 592 (2023).

\bibitem{EichhornLinoWagner2023}
A.~Eichhorn, R.~R.~Lino dos Santos, and F.~Wagner,
\textit{Shift-Symmetric Horndeski Gravity in the Asymptotic-Safety Paradigm},
JCAP \textbf{02}, 052 (2023).

\bibitem{EichhornGoldHeld2023}
A.~Eichhorn, R.~Gold, and A.~Held,
\textit{Horizonless Spacetimes As Seen by Present and Next-Generation
Event Horizon Telescope Arrays},
Astrophys.\ J.\ \textbf{950}, 117 (2023).

\bibitem{deBritoEichhornLino2022}
G.~P.~de~Brito, A.~Eichhorn, and R.~R.~Lino dos Santos,
\textit{Are There ALPs in the Asymptotically Safe Landscape?},
JHEP \textbf{06}, 013 (2022).

\bibitem{CarballoRubioEtal2025}
R.~Carballo-Rubio et al.,
\textit{Towards a Non-Singular Paradigm of Black Hole Physics},
JCAP \textbf{05}, 003 (2025).

%% --- Geometric Flows ---

\bibitem{Hamilton1982}
R.~S.~Hamilton,
\textit{Three Manifolds with Positive Ricci Curvature},
J.\ Differential Geom.\ \textbf{17}, 255--306 (1982).

\bibitem{Perelman2002}
G.~Perelman,
\textit{The Entropy Formula for the Ricci Flow and Its Geometric Applications},
arXiv:math/0211159 (2002).

%% --- Nonequilibrium Thermodynamics and Information Theory ---

\bibitem{Onsager1931}
L.~Onsager,
\textit{Reciprocal Relations in Irreversible Processes},
Phys.\ Rev.\ \textbf{37}, 405 (1931).

\bibitem{Prigogine1967}
I.~Prigogine,
\textit{Introduction to Thermodynamics of Irreversible Processes},
Interscience, New York (1967).

\bibitem{Jaynes1957}
E.~T.~Jaynes,
\textit{Information Theory and Statistical Mechanics},
Phys.\ Rev.\ \textbf{106}, 620 (1957).

\end{thebibliography}

\end{document}
